% ----------
% A LaTeX template for course project reports
%
% This template is modified from "Tech Report ala MIT AI Lab (1981)"
% ----------
\documentclass[12pt, letterpaper, oneside]{article}
\usepackage{geometry}
\usepackage[utf8]{inputenc}
\usepackage[english]{babel}
\usepackage[runin]{abstract}
\usepackage{titling}
\usepackage{booktabs}
\usepackage{fancyhdr}
\usepackage{helvet}
\usepackage{csquotes}
\usepackage{graphicx}

\usepackage{amsmath}
\usepackage{amssymb}
\usepackage{bm}
\usepackage{natbib}
\usepackage{parskip}
\usepackage{etoolbox}
\usepackage{xcolor}

% Mathematical Shortcuts
\newcommand{\yv}{\mathbf{y}}
\newcommand{\Rp}{R_{\!P}}
\newcommand{\logg}{\log g}
\newcommand{\Tint}{T_{\mathrm{int}}}
\newcommand{\Kzz}{\log K_{zz}}
\newcommand{\fhr}{f_{\mathrm{HR}}}
\newcommand{\XC}{\log X_{\mathrm{C}}}
\newcommand{\XO}{\log X_{\mathrm{O}}}
\newcommand{\XN}{\log X_{\mathrm{N}}}
\newcommand{\ph}[1]{#1}
\newcommand{\XS}{\log X_{\mathrm{S}}}

\input{preamble.tex}

% ----------
% Variables
% ----------

\title{\textbf{Multi-fidelity Gaussian
Process Emulation of Exoplanet Emission
Spectra}}
\runningtitle{MF GPs for Exoplanets}
\author{Camilla Andreozzi}
\runningauthor{Camilla Andreozzi}
\affiliation{ETH Zürich}
\department{D-MATH, Seminar for Statistics}
\theyear{2026}
\mydate{Sep 22, 2026}

% ----------
% Review apparatus
% ----------
% Every \reviewnote of the previous draft has been resolved against the code and
% the saved run outputs. \resolved records what was checked and where the
% evidence lives, so each claim in the text can be traced to a file. \openitem
% marks the few questions that the repository cannot answer; they concern
% external provenance, not the reported computation.

\newcommand{\resolved}[2]{\par\begingroup\small\color{green!35!black}%
\noindent\textbf{[RESOLVED: #1]}\par\nobreak\noindent #2\par\endgroup}
\newcommand{\openitem}[2]{\par\begingroup\small\color{red!65!black}%
\noindent\textbf{[OPEN: #1]}\par\nobreak\noindent #2\par\endgroup}

\begin{document}
\maketitle
\begin{abstract}
Atmospheric retrieval requires repeated evaluations of expensive forward
models. We investigate multi-fidelity Gaussian process surrogates for the
emission spectrum of WASP-80b, with nine atmospheric and planetary inputs
and 195 wavelength bins. The available data comprise 97 high-fidelity
(HF) spectra and a larger set of 10{,}000 low-fidelity (LF) spectra.
We compare per-wavelength autoregressive Gaussian processes with a model
that combines a mean shared between fidelities with an autoregressive
residual covariance. In the leave-one-out comparison, the boosted
multi-fidelity model has the lowest mean per-wavelength normalised
root-mean-square error, 0.4350, compared with 0.4864 for HF-only GPBoost
and 0.4826 for a PCA--PCE reference. Its mean absolute error is 86.9 ppm,
compared with a mean observational uncertainty of 91.1 ppm. The advantage
over the HF-only boosted baseline is not significant under a paired
design-level test ($p=0.42$), so the gain is carried by the per-wavelength
aggregation rather than by a uniform improvement across designs. Benefits
also vary across spectral bands. We formulate a retrieval likelihood that
carries emulator error, and report retrievals of a synthetic and the
observed spectrum; the implemented likelihood uses a full cross-wavelength
emulator error covariance estimated from leave-one-out residuals, because
the Gaussian process predictive variances are shown here to be
miscalibrated.
\end{abstract}

\resolved{reported evaluation}{Every number in the abstract and in
Section~\ref{sec:results} was recomputed from
\texttt{results/cv/loo\_lf10000/loo\_predictions.csv} and matches
\texttt{manuscript/cv\_metrics.csv} and
\texttt{manuscript/supporting\_statistics.json} exactly. All five table rows
come from that single prediction file, whose
\texttt{run\_config.json} records \texttt{lf\_sample\_size: 10000},
\texttt{n\_lf\_used: 10000}, \texttt{n\_hf\_used: 97},
\texttt{run\_model1: true}, \texttt{run\_model3: true}. The PCA--PCE row is
aligned in \texttt{recompute\_cv\_outputs.reference\_predictions}. Two claims
of the previous draft are corrected below: the surrogates do not use an SE-ARD
kernel, and the headline advantage over the HF-only boosted baseline is not
statistically significant. The held-out HF cases are excluded from the
spectral emulators, but not from the upstream LF profile emulators; see
Section~\ref{sec:cv_protocol}.}

\thispagestyle{firstpage}
\pagebreak
\newgeometry{}

\section{Introduction}
\label{sec:intro}

Atmospheric retrieval infers atmospheric and planetary parameters
$\bm\theta$ from an observed spectrum $\mathbf y_{\mathrm{obs}}$ using a
physical forward model $\mathbf f_H(\bm\theta)$. The target is
\[
p(\bm\theta\mid\mathbf y_{\mathrm{obs}})
\propto p(\mathbf y_{\mathrm{obs}}\mid\bm\theta)p(\bm\theta)
\]
\citep{BarstowHeng2020}. Different parameter combinations can produce
similar spectra, leading to parameter dependence and weak identifiability.
Posterior sampling requires repeated forward evaluations, which become
costly when the simulator enforces coupled atmospheric physics. A surrogate
approximates the parameter-to-spectrum map, but its accuracy is constrained
by the number of HF simulations available for training.

This study asks whether LF information improves prediction of complete HF
spectra at withheld atmospheric inputs. The application is WASP-80b, whose
emission spectrum is represented by 195 wavelength bins. The available
simulations comprise 97 HF cases and a larger LF design of 10{,}000 spectra.
The LF model uses approximated atmospheric profiles, reducing the cost of
the physical calculation while introducing spectral discrepancies. The
statistical task is to use this LF information and learn its relationship to
the HF response from the limited HF design.

\subsection{Contribution and outline}

We compare two per-wavelength formulations. Model~1 uses an autoregressive
Gaussian process (GP), expressing the HF response as a scaled LF process
plus an independent discrepancy \citep{KennedyOHagan2000}. Model~2 combines
a mean learned through tree boosting with autoregressive GP residuals,
following the combination of tree boosting and Gaussian processes in GPBoost
\citep{Sigrist2022gpboosting}. Benchmarks include HF-only GP and GPBoost
models and the PCA--PCE surrogate of \citet{Nussbaumer2026}.

The contribution is an empirical assessment of these formulations across
wavelengths and atmospheric inputs, including the scale of prediction error
relative to observational uncertainty, and an assessment of whether the
predictive variances these models report can be used as uncertainty in a
retrieval. The boosted MF model has the lowest mean per-wavelength
normalised RMSE, although the gains vary across bands and are not
significant against the HF-only boosted baseline at the design level. The
predictive variances are found to be miscalibrated, which determines how
the retrieval likelihood is built.

Section~\ref{sec:related work} positions the study relative to earlier
surrogates. Section~\ref{sec:background} describes the physical models,
datasets and exploratory analysis. Section~\ref{sec:surrogate} specifies
the surrogates as implemented and the evaluation protocol, and
Section~\ref{sec:results} presents the comparisons, including predictive
calibration. Section~\ref{sec:bayesian-inversion} describes their use in
retrieval. Section~\ref{sec:disc} discusses implications and limitations.
Appendices~\ref{app:ar_background} and \ref{subsec:model_augmented} provide
the scalar autoregressive background and a wavelength-augmented extension.

\section{Related Work}
\label{sec:related work}

\paragraph{Atmospheric surrogates.}
\label{sec:lit-surrogates}
Surrogates can approximate either the emergent spectrum or intermediate
atmospheric quantities. \citet{Himes2022} used a neural-network
radiative-transfer surrogate within Bayesian retrieval of HD~189733\,b,
obtaining marginal posteriors close to those of the original simulator.
\citet{Gebhard2024} learned a low-dimensional representation of physically
consistent pressure--temperature profiles. Grid interpolation offers another
approximation, but its design matters: \citet{FisherHeng2022} showed that
interpolation on rectilinear spectral grids can bias retrieved posteriors.

The closest application is \citet{Nussbaumer2026}, who investigated
PCA-based polynomial chaos expansion (PCE) and polynomial chaos--Kriging
(PCK) surrogates for WASP-80b. That work considered direct HF spectral
emulation and an LF atmospheric model with a learned spectral discrepancy;
the reported primary retrievals used the single-fidelity PCE surrogate.
It provides the reference for the present comparison of spectral
representations and fidelity-coupling models, and it also supplies the LF
data used here.

\paragraph{Fidelity coupling and spectral representation.}
\label{sec:lit-mf}
The autoregressive GP framework represents the HF response as a scaled LF
process plus an independent discrepancy \citep{KennedyOHagan2000}.
\citet{Forrester2007} applied co-kriging to multi-fidelity optimisation.
Under the required nesting and model assumptions, recursive co-kriging
reproduces the joint predictive mean and variance through smaller GP
problems \citep{LeGratiet2014,LeGratietGarnier2014}.
\citet{Perdikaris2017} extended this approach to nonlinear relationships
between fidelities. The nesting condition those recursive results rely on
does not hold for the designs used here (Section~\ref{sec:data}); the
models are therefore fitted as joint two-fidelity GPs rather than
recursively.

For spectral outputs, independent scalar models permit different behaviour
at each wavelength but omit cross-wavelength predictive covariance.
Basis approaches emulate a low-dimensional representation
\citep{Higdon2008}, while multi-output kernels model dependence between
responses \citep{Alvarez2012}. GPBoost provides a complementary extension
of the mean model: a boosted predictor and GP covariance parameters can be
estimated jointly \citep{Sigrist2022gpboosting}. The present boosted model
uses a tree mean with an autoregressive residual covariance, fitted in
stages as described in Section~\ref{subsec:model2}.

\paragraph{Uncertainty and retrieval.}
\label{sec:lit-validation}
Forward error can affect inference. Bounds connecting GP approximation
error to differences between exact and approximate Bayesian
posteriors \citep{StuartTeckentrup2018} motivate assessing uncertainty
alongside spectral accuracy. Section~\ref{sec:calibration} carries out that
assessment for the models reported here.

\openitem{bibliography}{The bibliography file is not part of the code
repository, so the citation keys cannot be checked here. Two items still need
a look before submission: whether \texttt{LeGratiet2014} and
\texttt{LeGratietGarnier2014} are distinct works or one duplicated entry, and
the exact JWST data product behind \texttt{data/Observed\_Spectra.csv}, which
carries no provenance header (see Section~\ref{sec:data}).}

\section{Physical Models and Data}
\label{sec:background}
\label{sec:data}

\begin{figure}
    \centering
    \includegraphics[width=0.95\linewidth]{figures/spectra_linear.png}
    \par\medskip
    \includegraphics[width=0.95\linewidth]{figures/spectra_log10.png}
    \caption{Simulated and observed emission spectra on the linear scale (top)
    and after a base-10 logarithmic transformation (bottom). Blue curves show
    the $97$ HF spectra. Orange curves show the matched LF evaluations and the
    larger LF design. Black points show the $195$ observed eclipse depths;
    their measurement uncertainties are not displayed. In the logarithmic
    panel only the $190$ positive observations are shown; five nonpositive
    values, the smallest $-71.4$ ppm, cannot be displayed on that scale and
    are omitted. Connecting lines guide the eye between sampled wavelengths.}
    \label{fig:spectral_overview}
\end{figure}

\subsection{Physical forward models}

The forward models map nine atmospheric and planetary parameters
$\bm\theta\in\mathcal{X}\subset\mathbb{R}^{9}$ to an emission spectrum. The
inputs describe planetary radius, surface gravity, internal temperature,
vertical mixing, heat redistribution, and the carbon-to-oxygen ratio together
with solar-relative nitrogen, oxygen and sulfur abundances; they are listed
with their ranges in Table~\ref{tab:inputs}. At fidelity $t\in\{L,H\}$, the
discretised forward map is
\[
    \mathbf{f}_t:\mathcal{X}\longrightarrow\mathbb{R}^{195},
    \qquad
    \mathbf{f}_t(\bm\theta)
    =\bigl(f_{t,1}(\bm\theta),\ldots,
           f_{t,195}(\bm\theta)\bigr)^\top.
\]
Each component is the eclipse depth at a specified wavelength bin. Eclipse
depths are dimensionless and are converted to parts per million (ppm) by
multiplication by $10^6$.

The high-fidelity (HF) simulator couples atmospheric chemistry and thermal
structure \citep{Nussbaumer2026}. FastChem provides initial
chemical-equilibrium abundances. HELIOS computes the radiative--convective
temperature profile using opacities from HELIOS-K. VULCAN models chemical
kinetics, photochemistry and vertical mixing. The temperature and abundance
profiles are iterated to consistency, then passed to petitRADTRANS to compute
the emission spectrum. The atmosphere is assumed cloud-free, with a
vertically constant eddy-diffusion coefficient.

The low-fidelity (LF) model replaces the iterative atmospheric calculation
with emulated temperature and molecular-abundance profiles. These profiles are
represented through principal components, whose coefficients are predicted by
polynomial chaos expansions trained on HF atmospheric profiles
\citep{Nussbaumer2026}. The reconstructed profiles are passed to
petitRADTRANS. The retained molecular species are
$\mathrm{H_2O}$, $\mathrm{CO}$, $\mathrm{CO_2}$, $\mathrm{CH_4}$,
$\mathrm{NH_3}$ and $\mathrm{SO_2}$. Both fidelities therefore use the same
radiative-transfer engine. Their differences arise from the
atmospheric-profile approximations, vertical discretisation and reduced
chemical representation.

\subsection{Input coordinates and design}

\begin{table}[t]
    \centering
    \small
    \caption{The nine input coordinates, as they appear in
    \texttt{data/XHF.csv} and \texttt{data/XLF10k.csv} (that column order),
    with the range spanned by the $97$ HF design points. These are the same
    coordinates in which the retrieval prior of
    Section~\ref{subsub:prior} is specified; no conversion is applied
    anywhere in the pipeline.}
    \label{tab:inputs}
    \begin{tabular}{llrr}
        \toprule
        Symbol & Meaning & HF min & HF max \\
        \midrule
        $\log_{10}K_{zz}$ & eddy diffusion, $\mathrm{cm^2\,s^{-1}}$ & 8.010 & 11.487 \\
        $R_P$             & planet radius, Jupiter radii            & 0.919 & 1.075 \\
        $T_{\mathrm{int}}$& internal temperature, K                 & 151.24 & 449.09 \\
        $\mathrm{C/O}$    & carbon-to-oxygen ratio                  & 0.123 & 1.498 \\
        $[\mathrm{N/H}]$  & solar-relative nitrogen, dex            & $-0.972$ & 1.974 \\
        $[\mathrm{O/H}]$  & solar-relative oxygen, dex              & $-0.993$ & 1.989 \\
        $[\mathrm{S/H}]$  & solar-relative sulfur, dex              & $-0.976$ & 1.999 \\
        $\log_{10}g$      & surface gravity, $\mathrm{cm\,s^{-2}}$  & 3.107 & 3.190 \\
        $f_{\mathrm{HR}}$ & heat redistribution factor              & 0.252 & 0.666 \\
        \bottomrule
    \end{tabular}
\end{table}

\resolved{inputs and provenance}{The nine coordinates were read from the
column headers of \texttt{data/XHF.csv} and \texttt{data/XLF10k.csv}, which
agree. The previous draft warned that the methods text might use elemental log
mixing ratios while the retrieval uses C/O and solar-relative abundances: that
mismatch does not exist in this pipeline. The training design is already in
C/O and $[\mathrm{X/H}]$ coordinates, identical to those of
\texttt{bayesian\_inversion/priors.py}, so no conversion is needed between
training and retrieval. The earlier elemental-mixing-ratio description
belonged to the thesis data generation, not to the files used here. The design
ranges above also coincide with the prior support of
Section~\ref{subsub:prior}, which is expected: the 10{,}000-point LF design is
itself a draw from that prior, a fact \texttt{priors.check\_box\_against\_design}
tests directly.}

\subsection{Simulation and observation data}

The available data comprise $97$ HF spectra, LF evaluations at those same $97$
inputs, and a separate LF design containing $10\,000$ spectra. All spectra are
represented on the same $195$-bin wavelength grid.

These three sets play different roles, and the distinction matters for the
model assumptions. The $97$ matched LF evaluations
(\texttt{data/YLF.csv}) are used only in the exploratory analysis of
Section~\ref{sec:eda}. Every fitted model in this paper draws its LF training
rows from the $10\,000$-point design (\texttt{data/XLF10k.csv},
\texttt{data/YLF10k.csv}); the matched evaluations never enter a fit.

\emph{The two designs are not nested.} Comparing the $97$ HF input vectors
with the $10\,000$ LF input vectors, the smallest Euclidean distance between
an HF point and its nearest LF point is $0.373$ in raw input units, with
median $1.011$; no HF input appears in the LF design even approximately. The
Markov and nesting conditions that make recursive co-kriging equivalent to the
joint model \citep{LeGratietGarnier2014} therefore do not hold, and the models
below are fitted as joint two-fidelity GPs on stacked LF and HF rows, which
does not require nesting. It also means that at every prediction point the LF
process is itself interpolated rather than observed.

The observed spectrum of WASP-80b combines JWST NIRCam F322W2, NIRCam F444W
and MIRI LRS measurements. It provides the eclipse-depth vector
$\mathbf{y}_{\mathrm{obs}}\in\mathbb{R}^{195}$ and its measurement
uncertainties for atmospheric retrieval. The file supplies asymmetric error
bars, \texttt{measured\_err\_lo} and \texttt{measured\_err\_hi}; throughout
this paper the per-bin uncertainty is their arithmetic mean,
$\sigma_{\mathrm{obs},j}=\tfrac12(\sigma^-_j+\sigma^+_j)$, and
$\bar\sigma_{\mathrm{obs}}=91.1$ ppm is the mean of that quantity over bins.
The surrogate models are fitted to simulator outputs; the astronomical
observations supply the retrieval target.

Table~\ref{tab:spectral_data} summarises the wavelength grid. The spacing
between bin centres is exactly $0.015\,\mu\mathrm{m}$ within both NIRCam
filters and exactly $0.25\,\mu\mathrm{m}$ within MIRI. The separation between
the last NIRCam and first MIRI centres is $0.2875\,\mu\mathrm{m}$. NIRCam
contributes $168$ of the $195$ bins ($86.2\%$). Consequently, averages that
weight bins equally assign most of their weight to NIRCam, despite MIRI's
wider wavelength coverage.

\begin{table}[t]
    \centering
    \small
    \caption{Wavelength sampling and HF--LF association.
    Wavelength ranges refer to bin centres. Correlations are
    calculated across the $97$ matched atmospheric inputs at each
    wavelength, then summarised by their within-band median.}
    \label{tab:spectral_data}
    \begin{tabular}{lrrrr}
        \toprule
        Instrument & Bins &
        Range ($\mu$m) &
        Spacing ($\mu$m) &
        Median $r_j$ \\
        \midrule
        NIRCam F322W2 & 103 & 2.4575--3.9875 & 0.015 & 0.809 \\
        NIRCam F444W  & 65  & 4.0025--4.9625 & 0.015 & 0.878 \\
        MIRI LRS      & 27  & 5.2500--11.7500 & 0.250 & 0.821 \\
        \bottomrule
    \end{tabular}
\end{table}

\resolved{LF counts and overlap}{All reported multi-fidelity results use
$n_L=10\,000$, the complete LF design. This was confirmed three ways:
\texttt{results/cv/loo\_lf10000/run\_config.json} records
\texttt{n\_lf\_used: 10000}; the reporting script
\texttt{modelling/cv/report\_cv\_performance.py} reads only
\texttt{results/cv/loo\_lf10000/loo\_predictions.csv}; and the worker processes
verify their LF sample against
\texttt{lf\_sample\_source\_indices.csv} before fitting
(\texttt{cv\_loo.verify\_worker\_lf\_sample}). The ``200 LF training cases''
and the \texttt{lf\_sample\_size: 1000} that the previous draft noticed belong
to other runs: the latter is in \texttt{results/cv/loo/run\_config.json},
written by a later Model~2 run into a directory whose Model~1 and Model~3
lanes were produced earlier at the full LF budget. That stale config is not
the configuration of any result reported here. On nesting: the $97$ paired LF
inputs do \emph{not} belong to the $10\,000$-input design, as quantified
above.}

\openitem{observation source}{\texttt{data/Observed\_Spectra.csv} contains
only four columns (wavelength, depth, lower and upper error) with no
instrument label, program identifier or reduction reference, and the
repository holds no accompanying provenance record. The exact data product and
the bin-response convention (point evaluation at the bin centre versus an
average over the instrument bandpass) must be supplied from outside the code
before submission; the latter also affects the wavelength-augmented extension
of Appendix~\ref{subsec:model_augmented}. What can be confirmed here is that
the file's $195$ wavelengths match the simulation grid to within $10^{-14}$
(asserted in \texttt{report\_cv\_performance.main}), that five depths are
nonpositive, and that the asymmetric bars are handled as stated above.}

For fidelity $t\in\{L,H\}$, write the data used in a particular fit as
\begin{equation}
\mathcal D=\mathcal D_L\cup\mathcal D_H,
\qquad
\mathcal D_t=\{(\bm\theta_i^t,\mathbf y_i^t)\}_{i=1}^{n_t},
\qquad \mathbf y_i^t\in\mathbb R^m,
\quad m=195,
\label{eq:inversion-training-data}
\end{equation}
and let $\mathbf X_t\in\mathbb R^{n_t\times d}$ collect its inputs,
with $d=9$ and design set
$\Theta_t=\{\bm\theta_i^t\}_{i=1}^{n_t}$.
In the leave-one-out evaluation $n_L=10\,000$ and $n_H=96$ in every fold; in
the final fits used for retrieval, $n_L=10\,000$ and $n_H=96$, one HF case
being reserved for a closure test (Section~\ref{sec:retrieval-results}).

\subsection{Exploratory analysis}
\label{sec:eda}

Figure~\ref{fig:spectral_overview} shows the simulated and observed spectra
across the common wavelength grid. The two fidelities share the broad spectral
structure, with wavelength-dependent differences in amplitude and spread.
Eclipse depths and their absolute dispersion generally increase towards longer
wavelengths. Pooled prediction errors in original units can therefore be
dominated by these wavelengths. Wavelength-normalised errors provide a
complementary assessment.

The logarithmic display makes relative variation at shorter wavelengths more
visible. Five observed eclipse depths are nonpositive and are omitted only from
this display. Their presence must be distinguished from missing observations
when considering transformations of the response.

To assess the relationship between fidelities independently of their common
wavelength trend, we calculate the Pearson correlation
across matched atmospheric inputs at each bin:
\[
    r_j
    =
    \operatorname{Corr}_{i=1,\ldots,97}
    \left(y^H_{i,j},y^L_{i,j}\right),
    \qquad j=1,\ldots,195.
\]
All correlations are positive, ranging from $0.711$ to $0.956$, with an
arithmetic mean of $0.842$ across wavelengths. NIRCam F444W has the highest
median correlation (Table~\ref{tab:spectral_data}). The wavelength-resolved
curves also show substantial variation within each instrument band
(Figure~\ref{fig:hf_lf_correlation}).

Correlation alone does not quantify spectral bias. Writing
$\Delta_{ij}=y^L_{ij}-y^H_{ij}$ for the discrepancy at the matched inputs, the
signed bias $\bar\Delta_j$ is positive on average, $+22.7$ ppm over all bins,
and ranges from $-184.7$ to $+399.8$ ppm across the spectrum. The discrepancy
root-mean-square, $[n_H^{-1}\sum_i\Delta_{ij}^2]^{1/2}$, averages $169.0$ ppm.
Both grow with wavelength: mean signed bias is $+10.3$, $+18.5$ and $+80.5$
ppm in F322W2, F444W and MIRI, and mean discrepancy RMS is $108.5$, $187.6$
and $354.9$ ppm. Relative to the HF spread at each wavelength the picture is
flatter, the discrepancy RMS averaging $0.66$, $0.57$ and $0.77$ of $s_j$ in
the three bands (Figure~\ref{fig:hf_lf_discrepancy}). The LF model is therefore biased, not merely noisy, and the
bias is a substantial fraction of the between-design variation the surrogate
must explain. This is what the autoregressive discrepancy term is asked to
absorb.

These associations support investigating LF information for HF prediction.
Correlation and bias alone do not determine the autoregressive scaling
coefficient. The predictive benefit of combining fidelities is therefore
assessed against HF-only models using held-out HF spectra.

\begin{figure}
    \centering
    \includegraphics[width=\linewidth]{figures/hf_lf_correlation.png}
    \caption{Pearson correlation between HF and LF eclipse depths across $97$
    matched atmospheric inputs at each wavelength. The upper panel shows the
    full spectral range. Lower panels show individual instrument bands on a
    common expanded correlation scale. Dashed horizontal lines denote band
    medians.}
    \label{fig:hf_lf_correlation}
\end{figure}

\begin{figure}
    \centering
    \includegraphics[width=\linewidth]{figures/hf_lf_discrepancy.png}
    \caption{Wavelength-resolved HF--LF discrepancy
    $\Delta_{ij}=y^L_{ij}-y^H_{ij}$ over the $97$ matched atmospheric inputs.
    Top: signed bias $\bar\Delta_j$, with the shaded ribbon spanning
    $\bar\Delta_j\pm$ the standard deviation of $\Delta_{ij}$ across
    atmospheres. Middle: discrepancy root-mean-square
    $[n_H^{-1}\sum_i\Delta_{ij}^2]^{1/2}$. Bottom: that RMS divided by the HF
    spread $s_j$ at the same wavelength. Dashed horizontal lines are band
    means; both bias and discrepancy RMS grow with wavelength while the
    relative panel is flatter.}
    \label{fig:hf_lf_discrepancy}
\end{figure}

\resolved{focused EDA diagnostic}{The requested wavelength-resolved signed
HF--LF bias and discrepancy RMSE are computed above from
\texttt{data/YHF.csv} and \texttt{data/YLF.csv}, the $97$ matched pairs, by
\texttt{eda/hf\_lf\_discrepancy.py}, which writes
Figure~\ref{fig:hf_lf_discrepancy} and the per-wavelength and per-instrument
tables to \texttt{results/eda/}. The
logarithmic-panel legend is corrected in the caption of
Figure~\ref{fig:spectral_overview} to distinguish $195$ available from $190$
positive observations.}

\section{Surrogate Modelling and Evaluation}
\label{sec:surrogate}

Both models act independently at each wavelength. The first places the
autoregressive structure directly on the responses; the second applies it
to residuals around a boosted mean. Generic scalar background is given in
Appendix~\ref{app:ar_background}. This section describes the models
\emph{as implemented and scored}; where the implementation departs from the
formulation of the previous draft, the departure is stated explicitly.

\subsection{Common covariance, approximation and preprocessing}
\label{subsec:common}

Both models are fitted with GPBoost 1.7.1.1 \citep{Sigrist2022gpboosting}
using its two-level autoregressive multi-fidelity covariance,
\texttt{ar1\_mf\_matern}, which implements exactly
$f_H(\bm\theta)=\rho f_L(\bm\theta)+\delta(\bm\theta)$ with $f_L$ and $\delta$
independent GPs sharing a base covariance family but carrying separate
parameter vectors, and with $\rho$ unrestricted in sign. The last column of
the GP coordinates is the fidelity indicator $\mathbf 1\{t=H\}$; it selects
covariance blocks and is not a continuous kernel coordinate. The HF-only
baselines use the same base covariance without the autoregressive wrapper.

The base covariance is an \emph{isotropic Matérn} function with smoothness
$\nu=3/2$ (\texttt{cov\_fct\_shape} $=1.5$),
\begin{equation}
    k(\bm\theta,\bm\theta')=\sigma^2
    \left(1+\frac{\sqrt3\,\|\bm\theta-\bm\theta'\|}{\ell}\right)
    \exp\!\left(-\frac{\sqrt3\,\|\bm\theta-\bm\theta'\|}{\ell}\right),
    \label{eq:matern32}
\end{equation}
with a single range $\ell$ shared by all nine input dimensions. Consequently
each wavelength carries six covariance parameters in the multi-fidelity case,
\begin{equation}
    \{\sigma^2_{L,j},\ell_{L,j},
      \sigma^2_{\delta,j},\ell_{\delta,j},\rho_j,
      \sigma^2_{\varepsilon,j}\},
    \label{eq:implemented-params}
\end{equation}
that is $195\times6=1170$ across the spectrum, and three
$\{\sigma^2_j,\ell_j,\sigma^2_{\varepsilon,j}\}$ in the single-fidelity case.

\resolved{kernel identity}{The previous draft specified squared-exponential
covariances with automatic relevance determination (SE-ARD) and
per-dimension length-scales $\bm\ell\in(0,\infty)^d$. Neither property is
present in the runs that produced the results. \texttt{src/model\_1.py} and
\texttt{src/model\_3.make\_model3\_gp\_model} both request
\texttt{cov\_function="ar1\_mf\_matern"} (\texttt{"matern"} for HF-only) with
\texttt{cov\_fct\_shape=1.5}, i.e. Matérn $\nu=3/2$, and the isotropic base,
not the available \texttt{ar1\_mf\_matern\_ard} variant. The saved parameter
table \texttt{results/cv/loo\_lf10000/loo\_fitted\_parameters.csv} confirms
this directly: its covariance columns are
\texttt{low\_GP\_var, low\_GP\_range, discrepancy\_GP\_var,
discrepancy\_GP\_range, rho} and a single \texttt{Error\_var} --- one range
per component, not nine. The SE-ARD description has been replaced by
\eqref{eq:matern32} throughout; Appendix~\ref{app:ar_background} keeps the
general form and notes which specialisation was used.}

\resolved{diagonal variance and noise structure}{There is one Gaussian noise
variance per wavelength, not the two fidelity-specific variances
$\sigma^2_{\varepsilon,L,j}$ and $\sigma^2_{\varepsilon,H,j}$ of the previous
draft: the fitted-parameter table carries a single \texttt{Error\_var} column
for every fit. It is an estimated nugget --- the Gaussian likelihood's error
variance, fitted by maximum likelihood --- covering both simulator-level
variability and model misfit; the simulators themselves are deterministic, so
it is not measurement noise. It is unrelated to the observational uncertainty
$\sigma_{\mathrm{obs},j}$. No separate numerical jitter is added by the
project code; any regularisation beyond this nugget is internal to GPBoost.
The equations below have been changed to a single $\sigma^2_{\varepsilon,j}$
accordingly. Fitted nuggets are small: the median \texttt{Error\_var} is
$5.1\times10^{-13}$ for Model~1 MF and $2.8\times10^{-12}$ for Model~2 MF, in
squared eclipse-depth units.}

\paragraph{Computational approximation (Vecchia).}
A Vecchia approximation is used, and must be reported: it replaces the exact
$n\times n$ Cholesky factorisation of the likelihood by an ordered conditional
factorisation
$p(\mathbf y)\approx\prod_{i}p(y_i\mid \mathbf y_{N(i)})$, where $N(i)$
indexes a bounded set of previously ordered neighbours. The settings are
GPBoost's, with the neighbour count set explicitly by the project code:
\texttt{gp\_approx="vecchia"}, \texttt{num\_neighbors=20}, and the package
default \texttt{vecchia\_ordering="random"}; prediction uses
\texttt{num\_neighbors\_pred} $=2\times20=40$ neighbours by package default.
Neighbours are selected by distance in the scaled GP coordinates, with the
fidelity indicator included in the coordinate vector for the multi-fidelity
covariance. The approximation applies to \emph{all} multi-fidelity fits, whose
$10\,096$ training rows make the dense factorisation expensive. The two models
differ for the single-fidelity baselines: Model~1's HF-only fits also use
Vecchia, whereas Model~2's HF-only fits use the exact Cholesky
(\texttt{gp\_approx="none"}), on the grounds recorded in
\texttt{src/model\_3.py} that at $96$ rows the two agreed to about five
significant figures while the exact factorisation was cheaper. The equations
in Sections~\ref{subsec:model1} and \ref{subsec:model2} describe the exact
model; the likelihood actually maximised is its Vecchia approximation.

\resolved{Vecchia}{The previous draft asked whether the approximation was used
and to remove the paragraph otherwise. It is used, by every reported
multi-fidelity run, so the paragraph is retained and completed from
\texttt{src/model\_1.fit\_model1}, \texttt{src/model\_3.make\_model3\_gp\_model}
and the GPBoost 1.7.1.1 signature. The ordering and prediction-neighbour
settings are package defaults; they are stated here as such rather than
inferred silently, and the one asymmetry between the two models' baselines is
flagged above because it makes the two single-fidelity rows of
Table~\ref{tab:cv_metrics} not strictly like-for-like.}

\paragraph{Preprocessing.} The responses are \emph{not} transformed: all models
are fitted to eclipse depths in their original units, and predictions and
predictive variances need no back-transformation. Only the GP input
coordinates are standardised, by centring and scaling each of the nine inputs
with the mean and population standard deviation ($\mathrm{ddof}=0$) of that
fit's \emph{training HF rows}, and applying that transform to both LF and HF
rows. Validation rows reuse the stored training statistics, so no held-out
information enters the transform. In Model~2 the standardisation is applied to
the residual GP coordinates only; the boosted trees and the tuning GP retain
raw inputs.

\resolved{preprocessing}{The previous draft stated that within each bin one
centring and scaling was shared by both fidelities on the \emph{response}, and
that predictive variances were multiplied by the squared response scale. No
response transformation exists in the code
(\texttt{src/model\_1.\_fit\_input\_scaler} and
\texttt{src/model\_3.fit\_model3} scale inputs only), so those sentences have
been removed. The scaler is fitted on training HF rows only and therefore
excludes the held-out HF case, as required.}

\subsection{Model 1: per-wavelength autoregressive GPs}
\label{subsec:model1}

Model~1 applies the scalar autoregressive construction of
Appendix~\ref{app:ar_background} independently at each wavelength bin. For
$j=1,\ldots,m$, let $f_{t,j}(\bm\theta)=[\mathbf{f}_t(\bm\theta)]_j$ denote
the response at wavelength $\lambda_j$ and fidelity $t\in\{L,H\}$, where
$\bm\theta\in\mathcal{X}\subset\mathbb{R}^{d}$, $d=9$, and $m=195$.
The model specifies
\begin{equation}
    f_{H,j}(\bm\theta)
    = \rho_j f_{L,j}(\bm\theta) + \delta_j(\bm\theta),
    \label{eq:model1_ar}
\end{equation}
with independent Gaussian process priors
\begin{equation}
    f_{L,j}\sim\mathcal{GP}(\mu_{L,j}(\bm\theta),k_{L,j}),
    \qquad
    \delta_j\sim\mathcal{GP}(\mu_{\delta,j}(\bm\theta),k_{\delta,j}),
    \qquad
    f_{L,j}\perp\delta_j.
    \label{eq:model1_gps}
\end{equation}
The mean functions are linear in the nine \emph{raw} (unstandardised) inputs,
with separate low- and high-fidelity coefficient blocks, so the process is not
zero-mean; the coefficients are profiled out with the covariance parameters by
the marginal likelihood. The coupling $\rho_j\in\mathbb{R}$ is specific to bin
$j$ and constant over $\bm\theta$. The kernels $k_{L,j}$ and $k_{\delta,j}$
are the isotropic Matérn-$3/2$ covariances \eqref{eq:matern32} on standardised
atmospheric inputs, with hyperparameters
$\bm\psi_{L,j}=\{\sigma^2_{L,j},\ell_{L,j}\}$ and
$\bm\psi_{\delta,j}=\{\sigma^2_{\delta,j},\ell_{\delta,j}\}$.
The bins are mutually independent conditional on their hyperparameters.
Thus, Model~1 consists of $m$ scalar two-fidelity GPs and does not model
covariance across wavelengths.

\resolved{evaluated Model 1}{The previous draft suspected a staged
implementation --- an LF GP layer, then a coupling estimated from paired
outputs, then separate discrepancy GPs. That is not what was scored.
\texttt{src/model\_1.fit\_model1} builds a single \texttt{gpb.GPModel} with the
\texttt{ar1\_mf\_matern} covariance over the stacked LF and HF rows and calls
\texttt{fit} once, so $\rho_j$, both covariance blocks and the noise variance
are estimated \emph{together} from one LF--HF marginal likelihood, exactly as
\eqref{eq:model1_lml} states. The joint description is therefore correct and
has been kept. Two corrections were still needed: the mean is linear and
fidelity-specific rather than zero (\texttt{fit\_model1} passes the raw inputs
as \texttt{X}, and leaves GPBoost's \texttt{fidelity\_specific\_mean=True}),
and there is one noise variance rather than two. Note that Model~1 and Model~2
differ here: Model~2 sets \texttt{fidelity\_specific\_mean=False}.}

At each wavelength, both models use
\begin{equation}
    \mathcal{D}_{t,j}
    = \{(\bm\theta_i^t,y_{i,j}^t)\}_{i=1}^{n_t},
    \qquad
    y_{i,j}^t=f_{t,j}(\bm\theta_i^t)+\varepsilon_{i,j},
    \qquad
    \varepsilon_{i,j}\sim
    \mathcal{N}(0,\sigma^2_{\varepsilon,j}),
    \label{eq:perbin-data}
\end{equation}
with noise terms independent of one another and of the latent processes, and
one noise variance shared by the two fidelities. Let
$\mathbf{X}_t\in\mathbb{R}^{n_t\times d}$ collect the atmospheric inputs, and
define
\begin{equation}
    \mathbf{y}_{(j)}^t
    = (y_{1,j}^t,\ldots,y_{n_t,j}^t)^\top,
    \qquad
    \mathbf{y}_{(j)}
    = \begin{bmatrix}
        \mathbf{y}_{(j)}^L \\
        \mathbf{y}_{(j)}^H
      \end{bmatrix}
    \in\mathbb{R}^{n_L+n_H}.
    \label{eq:model1-stacked}
\end{equation}

The joint formulation estimates the LF and discrepancy covariance parameters
together with $\rho_j$ through the LF--HF marginal likelihood. Writing
$\mathbf m_{(j)}$ for the stacked linear mean, the training distribution is
$\mathbf{y}_{(j)}\sim\mathcal{N}(\mathbf m_{(j)},
\bm\Sigma_j(\rho_j,\bm\eta_j))$, where
\begin{equation}
    \bm\Sigma_j =
    \begin{bmatrix}
        \mathbf{K}_{L,j}^{LL}
          + \sigma^2_{\varepsilon,j}\mathbf{I}_{n_L}
        & \rho_j\mathbf{K}_{L,j}^{LH} \\[4pt]
        \rho_j\mathbf{K}_{L,j}^{HL}
        & \rho_j^2\mathbf{K}_{L,j}^{HH}
          + \mathbf{K}_{\delta,j}^{HH}
          + \sigma^2_{\varepsilon,j}\mathbf{I}_{n_H}
    \end{bmatrix}.
    \label{eq:model1_blockcov}
\end{equation}
For $q\in\{L,\delta\}$ and $A,B\in\{L,H\}$,
$[\mathbf{K}_{q,j}^{AB}]_{ii'}
=k_{q,j}(\bm\theta_i^A,\bm\theta_{i'}^B)$.
The remaining hyperparameters are collected in
\begin{equation}
    \bm\eta_j
    = \{\bm\psi_{L,j},\bm\psi_{\delta,j},
         \sigma^2_{\varepsilon,j}\}.
    \label{eq:model1-hyperparameters}
\end{equation}
The log marginal likelihood is
\begin{equation}
    \log L_j(\rho_j,\bm\eta_j)
    = -\frac{1}{2}(\mathbf{y}_{(j)}-\mathbf m_{(j)})^\top
        \bm\Sigma_j^{-1}(\mathbf{y}_{(j)}-\mathbf m_{(j)})
      -\frac{1}{2}\log\det\bm\Sigma_j
      -\frac{n_L+n_H}{2}\log(2\pi),
    \label{eq:model1_lml}
\end{equation}
maximised in its Vecchia-approximated form. Each bin is fitted independently by
\begin{equation}
    (\widehat\rho_j,\widehat{\bm\eta}_j)
    = \arg\max_{\rho_j,\bm\eta_j}
      \log L_j(\rho_j,\bm\eta_j),
    \qquad j=1,\ldots,m.
    \label{eq:model1_criterion}
\end{equation}
Optimisation is GPBoost's gradient-based default. Starting values are set from
the data scale by the project code
(\texttt{\_initial\_cov\_pars\_for\_gp\_model} in Model~2; Model~1 uses
GPBoost's own defaults unless warm-started, and warm starting was disabled,
\texttt{warm\_start\_model1\_mf: false}): variances are initialised from the
response variance, ranges from a median inter-point distance, and $\rho$ at
$1$.

For prediction at $\bm\theta^*$, define the training-to-HF covariance vector
\begin{equation}
    \mathbf{c}_{H,j}(\bm\theta^*)
    = \begin{bmatrix}
        \rho_j k_{L,j}(\mathbf{X}_L,\bm\theta^*) \\
        \rho_j^2 k_{L,j}(\mathbf{X}_H,\bm\theta^*)
          + k_{\delta,j}(\mathbf{X}_H,\bm\theta^*)
      \end{bmatrix}.
    \label{eq:model1-cross}
\end{equation}
Gaussian conditioning gives
\begin{equation}
    f_{H,j}(\bm\theta^*)\mid
    \mathcal{D}_{L,j},\mathcal{D}_{H,j},
    \widehat\rho_j,\widehat{\bm\eta}_j
    \sim\mathcal{N}\bigl(
        \mu_{H,j}(\bm\theta^*),s^2_{H,j}(\bm\theta^*)\bigr),
    \label{eq:model1-pred}
\end{equation}
with
\begin{align}
    \mu_{H,j}(\bm\theta^*)
    &= m_{H,j}(\bm\theta^*)+\mathbf{c}_{H,j}(\bm\theta^*)^\top
       \bm\Sigma_j^{-1}(\mathbf{y}_{(j)}-\mathbf m_{(j)}),
       \label{eq:model1-mean} \\
    s^2_{H,j}(\bm\theta^*)
    &= \rho_j^2 k_{L,j}(\bm\theta^*,\bm\theta^*)
       + k_{\delta,j}(\bm\theta^*,\bm\theta^*) \nonumber\\
    &\quad - \mathbf{c}_{H,j}(\bm\theta^*)^\top
       \bm\Sigma_j^{-1}\mathbf{c}_{H,j}(\bm\theta^*).
       \label{eq:model1-var}
\end{align}
The saved predictions use GPBoost's \texttt{predict\_response=True}, so the
recorded variance includes $\sigma^2_{\varepsilon,j}$; the calibration
assessment of Section~\ref{sec:calibration} uses those recorded values.

\subsection{Model 2: boosted mean and autoregressive residuals}
\label{subsec:model2}

Model~2 combines a mean learned by tree boosting with autoregressive Gaussian
process residuals. As in Model~1, the $m=195$ wavelength bins are modelled
independently. For each bin $j$, we specify
\begin{align}
    f_{L,j}(\bm\theta)
    &= F_j(\bm\theta)+u_j(\bm\theta),
       \label{eq:model2-lf} \\
    f_{H,j}(\bm\theta)
    &= F_j(\bm\theta)+\rho_j u_j(\bm\theta)
       +\delta_j(\bm\theta),
       \label{eq:model2-hf}
\end{align}
where $F_j:\mathcal{X}\to\mathbb{R}$ is one mean function shared by the two
fidelities at bin $j$, and
\begin{equation}
    u_j\sim\mathcal{GP}(0,k_{L,j}),
    \qquad
    \delta_j\sim\mathcal{GP}(0,k_{\delta,j}),
    \qquad
    u_j\perp\delta_j.
    \label{eq:model2-gps}
\end{equation}
The process $u_j=f_{L,j}-F_j$ is the LF residual, while $\delta_j$
is an additional HF residual independent of $u_j$.
Conditional on $F_j$, both fidelities have mean $F_j(\bm\theta)$.
Equivalently,
\begin{equation}
    f_{H,j}(\bm\theta)
    = \rho_j f_{L,j}(\bm\theta)
      +(1-\rho_j)F_j(\bm\theta)+\delta_j(\bm\theta).
    \label{eq:model2-ar}
\end{equation}
Setting $F_j\equiv0$ recovers the Model~1 covariance structure. The coupling
$\rho_j$, kernels $k_{L,j}$, $k_{\delta,j}$ and hyperparameters $\bm\eta_j$
follow Section~\ref{subsec:common}. The shared-mean property is imposed
explicitly in code by \texttt{fidelity\_specific\_mean=False}.

Let $\mathcal{H}=\operatorname{span}(\mathcal{T})$, where $\mathcal{T}$ is the
class of regression trees \citep{Breiman1984}, and take $F_j\in\mathcal{H}$.
For a candidate mean function, define
\begin{equation}
    \mathbf{F}_{(j)}
    = \begin{bmatrix}
        F_j(\mathbf{X}_L) \\
        F_j(\mathbf{X}_H)
      \end{bmatrix},
    \qquad
    \mathbf{r}_{(j)}(F_j)
    = \mathbf{y}_{(j)}-\mathbf{F}_{(j)}.
    \label{eq:model2-residuals}
\end{equation}
Integrating out $u_j$ and $\delta_j$ gives
\begin{equation}
    \mathbf{y}_{(j)}\mid F_j,\rho_j,\bm\eta_j
    \sim\mathcal{N}\bigl(
        \mathbf{F}_{(j)},\bm\Sigma_j(\rho_j,\bm\eta_j)\bigr),
    \label{eq:model2-training}
\end{equation}
with the covariance in \eqref{eq:model1_blockcov}, and log marginal likelihood
\begin{equation}
    \log L_j(F_j,\rho_j,\bm\eta_j)
    = -\frac{1}{2}\mathbf{r}_{(j)}(F_j)^\top
        \bm\Sigma_j^{-1}\mathbf{r}_{(j)}(F_j)
      -\frac{1}{2}\log\det\bm\Sigma_j
      -\frac{n_L+n_H}{2}\log(2\pi).
    \label{eq:model2-lml}
\end{equation}

\paragraph{Estimation as implemented.} The fitted procedure is staged rather
than alternating, and the mean is learned from HF rows only:
\begin{enumerate}
\itemsep2pt
\item \emph{Tune.} Boosting hyperparameters are selected by Optuna TPE with
$100$ trials, scoring $5$-fold cross-validated mean squared error on the $96$
training HF rows, with up to $1000$ boosting rounds and early stopping after
$20$ rounds without improvement. Each trial is scored with a GP in the loop
(a single-fidelity Matérn GP on the HF rows, refitting its covariance
parameters within the trial), so the tuner does not select a tree that absorbs
the smooth structure the GP exists to model. The search ranges are GPBoost's
template ranges with the learning rate capped at $1$; the tuned quantities are
the learning rate, minimum data per leaf, number of leaves, $L_2$
regularisation, maximum bin count and feature fraction, with
\texttt{max\_depth} unrestricted. Tuning is repeated independently in every
fold and at every wavelength, and one tuning result is shared by the
single- and multi-fidelity variants of that fold and wavelength.
\item \emph{Fit the mean.} The tree ensemble $\widehat F_j$ is trained by
ordinary $L_2$ boosting on the training HF rows alone, with no GP in the
objective.
\item \emph{Fit the covariance.} $\widehat F_j$ is evaluated at all LF and HF
training inputs and supplied as a fixed offset; the autoregressive covariance
parameters $(\rho_j,\bm\eta_j)$ are then estimated by maximising
\eqref{eq:model2-lml} in its Vecchia-approximated form with
$\mathbf{F}_{(j)}$ held fixed.
\end{enumerate}
The estimates $(\widehat F_j,\widehat\rho_j,\widehat{\bm\eta}_j)$ are
therefore a one-pass staged solution of \eqref{eq:model2-lml}, not a
stationary point of the joint objective.

\resolved{evaluated Model 2}{The previous draft described alternating joint
estimation --- covariance update, then functional boosting update, iterated.
\texttt{src/model\_3.fit\_model3} does not do this: it calls \texttt{gpb.train}
without a \texttt{gp\_model} argument, so the booster is fitted by plain $L_2$
boosting, and only afterwards is \texttt{gp\_model.fit} called with the tree
predictions as \texttt{offset}. There is no outer loop. The joint-estimation
equations have been replaced by the three-stage description above. Two further
properties were verified rather than assumed. The mean is shared between
fidelities, as the draft claimed --- \texttt{make\_model3\_gp\_model} sets
\texttt{fidelity\_specific\_mean=False} --- but it is \emph{trained} on HF rows
only (\texttt{\_prepare\_model3\_training\_arrays} returns
\texttt{x\_tree}=HF rows, \texttt{x\_gp\_offset}=all rows), so LF data enter
Model~2 exclusively through the covariance. The residual covariance is the
same \texttt{ar1\_mf\_matern} used by Model~1. Tuning settings are the
defaults of \texttt{tune\_model3\_parameters}, since the run configuration
records \texttt{model3\_tuning\_kwargs: null} and
\texttt{model3\_gp\_kwargs: null}.}

Prediction uses \eqref{eq:model1-cross} evaluated at the Model~2 parameters.
With $\widehat{\mathbf{r}}_{(j)}=\mathbf{r}_{(j)}(\widehat F_j)$,
\begin{align}
    \mu_{H,j}(\bm\theta^*)
    &= \widehat F_j(\bm\theta^*)
       +\mathbf{c}_{H,j}(\bm\theta^*)^\top
        \bm\Sigma_j^{-1}\widehat{\mathbf{r}}_{(j)},
       \label{eq:model2-mean} \\
    s^2_{H,j}(\bm\theta^*)
    &= \rho_j^2 k_{L,j}(\bm\theta^*,\bm\theta^*)
       +k_{\delta,j}(\bm\theta^*,\bm\theta^*)
       -\mathbf{c}_{H,j}(\bm\theta^*)^\top
       \bm\Sigma_j^{-1}\mathbf{c}_{H,j}(\bm\theta^*).
       \label{eq:model2-var}
\end{align}
The full spectrum is assembled bin by bin.

Both models condition on estimated covariance parameters. Model~2 also
conditions on $\widehat F_j$, so its predictive variance does not propagate
uncertainty in the fitted tree ensemble, and because the mean is fitted before
the covariance, the covariance parameters absorb no uncertainty about it
either. Neither predictive variance integrates over hyperparameter
uncertainty. Section~\ref{sec:calibration} measures the consequence.

\subsection{Implementation and HF-only benchmarks}

The comparison includes a per-wavelength HF-only GP, a per-wavelength
HF-only GPBoost model, and the PCA--PCE reference of
\citet{Nussbaumer2026}. These baselines distinguish improvements associated
with LF information from those associated with the mean model or spectral
representation. Table~\ref{tab:run_spec} collects the settings behind each
reported row.

\begin{table}[t]
  \centering
  \small
  \setlength{\tabcolsep}{4pt}
  \caption{Settings behind each row of Table~\ref{tab:cv_metrics}. All rows
  are scored on the same $97$ leave-one-out folds and the same $195$
  wavelengths, from
  \texttt{results/cv/loo\_lf10000/loo\_predictions.csv}. Software: GPBoost
  1.7.1.1, Optuna $\ge$3.4, SciPy $\ge$1.11, NumPy $\ge$1.26.}
  \label{tab:run_spec}
  \begin{tabular}{llllll}
    \toprule
    Row & $n_L$ & $n_H$ & Mean model & Covariance & Approximation \\
    \midrule
    per-$\lambda$ GP        & 0      & 96 & linear (raw inputs) & Matérn 3/2 & Vecchia, 20 \\
    per-$\lambda$ GPBoost   & 0      & 96 & boosted trees, HF   & Matérn 3/2 & exact \\
    PCA--PCE$^{\dagger}$    & ---    & 96 & PCE on PCA scores   & ---        & --- \\
    Model 1 MFGP            & 10\,000 & 96 & linear, fidelity-specific & AR(1) Matérn 3/2 & Vecchia, 20 \\
    Model 2 GPBoost         & 10\,000 & 96 & boosted trees, HF, shared & AR(1) Matérn 3/2 & Vecchia, 20 \\
    \bottomrule
  \end{tabular}

  \vspace{3pt}
  \begin{minipage}{0.95\textwidth}
    \footnotesize
    Inputs standardised on training HF rows; responses untransformed.
    Boosting tuned per fold and wavelength by TPE, $100$ trials, $5$-fold CV
    on MSE, $\le1000$ rounds, early stopping $20$. Covariance parameters,
    couplings, boosting hyperparameters and input scalers are all refitted
    within each fold. $^{\dagger}$Supplied predictions, not refitted here.
  \end{minipage}
\end{table}

\resolved{run specification}{Table~\ref{tab:run_spec} answers the request for
a compact settings table. Two points on how settings were chosen without
held-out targets: boosting hyperparameters are selected by cross-validation
\emph{within} each fold's $96$ training HF rows, so the held-out case is not
used; and the GP covariance parameters are maximum-likelihood estimates on the
training rows, with no tuning against held-out error. Case and wavelength
alignment is asserted in code rather than assumed ---
\texttt{report\_cv\_performance} checks that each predictor's recovered
\texttt{y\_true} matrix equals \texttt{data/YHF.csv} to $10^{-14}$, and
\texttt{reference\_predictions} checks the reference file's wavelength headers
against the HF grid. The PCA--PCE row is a supplied prediction file whose
internal configuration is not reproducible here; its row ordering remains an
assumption, discussed below.}

\subsection{Cross-validation protocol and metrics}
\label{sec:cv_protocol}

The evaluation leaves out complete HF spectra. Let
$\widehat f^{(-i)}_{H,j}(\bm\theta_i^H)$ denote the prediction at bin $j$
for HF case $i$ from a fit excluding that HF spectrum, and define
\[
e_{ij}=\widehat f^{(-i)}_{H,j}(\bm\theta_i^H)-y^H_{i,j},
\qquad i=1,\ldots,n_H,\quad j=1,\ldots,m,
\]
with $n_H=97$ reference cases and $m=195$ bins.

\paragraph{What is refitted in each fold.} One HF spectrum is removed; all
$10\,000$ LF rows are retained in every fold. Within a fold, and independently
at each wavelength, the following are refitted from the retained data alone:
the input scaler (from the $96$ training HF rows), the boosting
hyperparameters (by inner $5$-fold CV on those rows), the tree ensemble, the
covariance parameters $\bm\eta_j$, and the coupling $\rho_j$. Nothing is
carried across folds; warm starting was disabled.

\paragraph{What is not refitted.} The LF data themselves are a fixed input.
The LF spectra were generated by profile emulators that \citet{Nussbaumer2026}
trained on HF atmospheric profiles, and those emulators are not rebuilt here
--- the repository contains their outputs, not the generator. If the held-out
HF case contributed to training those upstream emulators, then the evaluation
measures generalisation \emph{conditional on that pretrained LF resource},
not full-pipeline generalisation to an HF case unseen by every learned
component. Two observations bound how much this matters. First, the $10\,000$
LF design contains no HF input, even approximately
(Section~\ref{sec:data}), so no LF row is an emulated copy of a held-out HF
spectrum; any leakage is indirect, through the profile emulators' fitted
coefficients. Second, the single-fidelity rows of Table~\ref{tab:cv_metrics}
use no LF data at all and are unaffected, which makes the MF-versus-SF
comparisons within this table the safest reading of the results.

\openitem{upstream LF provenance}{Which HF cases trained the LF profile
emulators cannot be determined from this repository; the thesis describes
emulators trained on all $97$. Until that is confirmed, the multi-fidelity
rows should be described as conditional on the pretrained LF resource, as
above, and the paper should not claim full-pipeline generalisation.}

For each wavelength, let
\[
\bar y^H_j=\frac1{n_H}\sum_{i=1}^{n_H}y^H_{i,j},
\qquad
s_j=\left[\frac1{n_H}\sum_{i=1}^{n_H}
(y^H_{i,j}-\bar y^H_j)^2\right]^{1/2}.
\]
The primary metric is the arithmetic mean of per-wavelength normalised
root-mean-square errors:
\begin{equation}
\mathrm{NRMSE}_j=
\frac{\left(n_H^{-1}\sum_{i=1}^{n_H}e_{ij}^2\right)^{1/2}}{s_j},
\qquad
\overline{\mathrm{NRMSE}}_{\mathcal J}
=\frac1{|\mathcal J|}\sum_{j\in\mathcal J}\mathrm{NRMSE}_j.
\label{eq:NRMSE}
\end{equation}
Here $\mathcal J$ contains all bins or those in an instrument band. The
population standard deviation ($\mathrm{ddof}=0$) is computed from the
reference values at each wavelength. This descriptive scoring convention
is separate from any scaling used to fit a model. It is neither a pooled
RMSE divided by the global response range nor a pooled RMSE divided by the
standard deviation of all spectral values; that pooled quantity is also
computed and stored, as \texttt{global\_nrmse} in
\texttt{manuscript/cv\_metrics.csv}, and is not what
Table~\ref{tab:cv_metrics} reports. No $s_j$ is zero: the smallest is
$3.8\times10^{-5}$ in eclipse-depth units, and no HF spectrum has zero norm
(smallest $\|\mathbf y^H_i\|_2=1.2\times10^{-2}$), so no degenerate-denominator
handling is needed.

Mean absolute error is reported in ppm:
\begin{equation}
\mathrm{MAE}_{\mathcal J}=
\frac1{n_H|\mathcal J|}\sum_{i=1}^{n_H}\sum_{j\in\mathcal J}|e_{ij}|.
\label{eq:MAE}
\end{equation}
Equal-bin averages assign 168 of 195 bins to NIRCam and do not weight by
wavelength interval. The reported ratio
$\mathrm{MAE}/\bar\sigma_{\mathrm{obs}}$ uses the mean measurement
uncertainty, $\bar\sigma_{\mathrm{obs}}=91.1$ ppm. It compares error scales;
it is not a predictive variance or a calibration diagnostic. Calibration is
assessed separately in Section~\ref{sec:calibration}.

The design-level comparison uses
\[
q_i=\left[\frac1m\sum_{j=1}^{m}(e_{ij}/s_j)^2\right]^{1/2}.
\]
This score is distinct from \eqref{eq:NRMSE} and from normalising a
spectrum's RMSE by its own standard deviation. The two-sided Wilcoxon
signed-rank tests compare 97 paired design scores, with Holm adjustment over
the comparisons reported in Section~\ref{sec:cv_results}. Standard deviations
of paired differences use denominator $n_H-1$. These tests are exploratory
because the LOO training sets overlap; pairing does not make the fitted models
independent across folds.

The atmospheric-input diagnostic uses the relative spectral error
\[
\frac{\|\widehat{\mathbf y}^{(-i)}_i-\mathbf y^H_i\|_2}
{\|\mathbf y^H_i\|_2},
\]
which has a different denominator and interpretation from both preceding
standardised scores.

\resolved{complete fold construction}{Answered by the two paragraphs above,
read from \texttt{modelling/cv/cv\_loo.py} (\texttt{run\_model1\_lane},
\texttt{run\_model3\_lane}, \texttt{split\_frame}) and
\texttt{src/cv.iter\_loo\_splits}. The refitting question and the upstream-LF
question have different answers, and the text now separates them.}

\resolved{central paired comparison}{The comparison the previous draft asked
for --- MF GPBoost against HF-only GPBoost --- was missing from the four
preregistered pairs and is now computed and reported in
Section~\ref{sec:cv_results}. It changes the reading of the headline result.
To keep the Holm family well defined, it is reported as a separate,
explicitly labelled comparison rather than folded into the original four;
the four original $p$-values are unchanged.}

\section{Results}
\label{sec:results}

\subsection{Cross-validated surrogate performance}
\label{sec:cv_results}

Table~\ref{tab:cv_metrics} reports the leave-one-out comparison over 97 HF
cases, using the definitions in Section~\ref{sec:cv_protocol}.

Under mean per-wavelength NRMSE, the ordering is Model~2 ($0.4350$),
the PCA--PCE reference of \citet{Nussbaumer2026} ($0.4826$),
Model~1 ($0.4837$), the single-fidelity GPBoost baseline
($0.4864$), and the single-fidelity GP baseline
($0.4953$). Model~2 improves on the reference by $9.9\%$.

\resolved{result-row identity}{The row labelled Model~2 in this paper is the
run stored as \texttt{model3} with \texttt{variant=mf}. The naming is a
historical artefact of the code: \texttt{results/cv/loo\_lf10000/run\_config.json}
records \texttt{run\_model3: true} and \texttt{run\_model2: false}, so the
repository's ``Model 2'' (a separate wavelength-augmented experiment) did not
run at all in this evaluation, and the boosted formulation of
Section~\ref{subsec:model2} is the repository's \texttt{model3}. The formulation
described there matches that code, with the staging correction recorded above.}

\subsection{Paired comparisons across atmospheric inputs}

Model~2 improves on Model~1 on $66$ of $97$ designs, with a mean
paired difference $q_i^{(1)}-q_i^{(2)}$ of $0.0706$
(sd $0.1701$; unadjusted $p=5.97\times10^{-5}$,
adjusted $p=2.39\times10^{-4}$). Model~1 improves on the PCA--PCE
reference on only $36$ designs: the mean paired difference
favours the reference by $0.0519$
(unadjusted $p=0.0045$, adjusted $p=0.0135$).
Model~2 improves on the reference on $55$ designs
(unadjusted and adjusted $p=0.2329$), and on the
single-fidelity GP on $56$ designs
(unadjusted $p=0.0224$, adjusted $p=0.0449$).

Two further comparisons, computed here and reported outside the Holm family
above, qualify the headline result. \emph{Model~2 does not significantly
improve on the HF-only GPBoost baseline}: it wins on $52$ of $97$ designs,
mean paired difference $0.0169$ (sd $0.1522$), Wilcoxon $p=0.419$, paired $t$
$p=0.276$. The same holds under MAE, where Model~2 wins on $51$ of $97$
designs with a mean advantage of $6.5$ ppm ($p=0.143$). Likewise Model~1 MF
does not significantly improve on its own single-fidelity counterpart:
it wins on $38$ of $97$ designs, Wilcoxon $p=0.121$.

The pattern is therefore that the boosted mean, not the fidelity coupling,
carries the improvement. Model~2 beats Model~1 decisively, and both boosted
rows beat the plain GP; but neither multi-fidelity model separates from its
own single-fidelity baseline at the design level. The $0.0514$ gap in mean
per-wavelength NRMSE between Model~2 and HF-only GPBoost arises from averaging
$\mathrm{NRMSE}_j$ over wavelengths, which weights bins where the MF model
happens to do well, and does not survive aggregation to one score per design.
All of these are exploratory comparisons under the overlapping-fold protocol
of Section~\ref{sec:cv_protocol}.

\subsection{Spectral variation and observational error scale}

MAE gives a different ordering among the remaining architectures. Model~2
leads at $86.9$ ppm, whereas Model~1 reaches $100.5$ ppm and is thereby worse
than the single-fidelity GP baseline ($99.8$ ppm).
The band-specific NRMSE minima also differ: Model~2 performs best in
F322W2, Model~1 in F444W, and the HF-only GP in MIRI
(Table~\ref{tab:cv_metrics}).
For Model~2, mean per-wavelength NRMSE is similar in the NIRCam and MIRI bands
($0.4352$ and $0.4339$), while MAE increases from $24.7$ ppm for bin centres
below $3\,\mu$m to $232.4$ ppm for bin centres strictly above $10\,\mu$m.
Over those same slices, the mean high-fidelity eclipse depth increases from
$171.1$ to $3832.0$ ppm, a factor of $22.4$. Long-wavelength bins therefore
account for much of the absolute error without having larger band-averaged
standardised error.

The best surrogate reaches $\mathrm{MAE}=86.9$ ppm, or
$0.954\,\bar\sigma_{\mathrm{obs}}$. Its per-bin MAE is below that bin's
observational uncertainty in $117$ of $195$ bins; its per-bin RMSE is below
that uncertainty in $75$ bins. Emulation error is therefore comparable to the
observational error scale, which is why the retrieval likelihood of
Section~\ref{sec:bayesian-inversion} carries an emulator error term.

For Model~2, squared error is concentrated in a subset of wavelength bins: the
worst $10$ bins account for $23.5\%$ of the total, the worst $5$ for $14.7\%$,
and the worst $20$ for $36.4\%$. The five largest squared-error contributions
occur between $10.5$ and $11.75\,\mu$m. The top ten also include one
shorter-wavelength bin, at $3.9875\,\mu$m, the last bin of F322W2. The largest
standardised error is $\mathrm{NRMSE}_j=0.6643$, at $4.0625\,\mu$m, the third
bin of F444W. Absolute-error concentration and standardised-error maxima
therefore identify different parts of the spectrum, and both maxima sit at or
near an instrument boundary.

\subsection{Predictive calibration}
\label{sec:calibration}

The retrieval construction of Section~\ref{sec:bayesian-inversion} uses
predictive variances, so those variances are assessed directly. Using the
saved leave-one-out predictive variances, define the standardised residual
$z_{ij}=e_{ij}/\widehat\sigma_{H,j}(\bm\theta_i^H)$ over all
$97\times195=18\,915$ held-out bin predictions. A calibrated Gaussian
predictive distribution would give $\mathbb E[z^2]=1$ and nominal interval
coverage.

\begin{table}[t]
  \centering
  \small
  \caption{Calibration of the leave-one-out predictive distributions, over all
  $18\,915$ held-out bin predictions. Coverage is the fraction of $|z_{ij}|$
  below the standard normal quantile. All four models are anticonservative in
  the tails; the two boosted models additionally produce a small number of
  near-zero variances, which inflates $\overline{z^2}$ without affecting the
  quantiles.}
  \label{tab:calibration}
  \begin{tabular}{lrrrrr}
    \toprule
    Model & $\overline{z^2}$ & $\mathrm{med}(z^2)$
          & 68\% & 95\% & 99\% \\
    \midrule
    per-$\lambda$ GP (SF)      & $1.28$              & 0.380 & 0.713 & 0.924 & 0.966 \\
    per-$\lambda$ GPBoost (SF) & $2.6\times10^{11}$  & 0.356 & 0.703 & 0.891 & 0.926 \\
    Model 1 MFGP               & $3.48$              & 1.014 & 0.498 & 0.778 & 0.874 \\
    Model 2 GPBoost            & $3.9\times10^{7}$   & 0.506 & 0.629 & 0.865 & 0.918 \\
    \bottomrule
  \end{tabular}
\end{table}

None of the four predictive distributions is calibrated
(Table~\ref{tab:calibration}). Every model under-covers at the $95\%$ and
$99\%$ levels, by between $3$ and $17$ percentage points, so the predictive
intervals are too narrow in the tails, which is the direction that matters for
a likelihood. Model~1 MF is the closest to nominal in the centre
($\mathrm{med}(z^2)=1.01$) but the worst in the tails ($95\%$ coverage
$0.778$): its errors are heavy-tailed relative to its own predictive scale.
The single-fidelity GP is the best calibrated overall and is still
anticonservative. The two boosted models produce occasional near-degenerate
variances --- $0.8\%$ of bins have a predictive standard deviation below $1\%$
of $s_j$, the smallest variance being $10^{-23}$ --- which contributes
essentially all of their enormous $\overline{z^2}$; their median and quantile
behaviour is unremarkable by comparison. This is consistent with
Section~\ref{subsec:model2}: the boosted predictive variance conditions on a
fitted tree ensemble whose uncertainty it does not carry, so at inputs the
trees fit closely the residual GP can report near-zero uncertainty.

The practical consequence is that $\bm\Sigma_{\mathrm{sur}}(\bm\theta)$ as
defined by \eqref{eq:model1-var} or \eqref{eq:model2-var} should not be used
as the emulator error term in a retrieval likelihood. The implemented
retrieval therefore replaces it by an empirical covariance of leave-one-out
residuals, as described in Section~\ref{subsec:implemented-likelihood}.

\subsection{Difficult atmospheric inputs}

Across design points, Model~2's spectrum-level relative $L_2$ error has median
$0.0923$, mean $0.1002$, empirical skewness $+4.11$, and maximum $0.492$. This
is a separate metric from standard-deviation-normalised RMSE.
Relative $L_2$ error correlates with oxygen abundance (Spearman
$\rho=+0.432$, $p\simeq 1.0\times10^{-5}$); the other eight inputs have
$|\rho|<0.08$ and $p>0.48$. The eight worst designs have mean oxygen
abundance $+1.19$ population standard deviations above the design mean. These
diagnostics associate larger errors with the high-$[\mathrm{O/H}]$ part of the
design domain, without establishing a causal explanation or demonstrating that
errors are confined to that region.

\begin{table}
  \centering
  \small
  \setlength{\tabcolsep}{5pt}
  \caption{Leave-one-out cross-validation errors over $n_H=97$
  high-fidelity designs. The NRMSE columns report arithmetic means of
  per-wavelength $\mathrm{RMSE}_j/s_j$, using the population standard
  deviation $s_j$ of the held-out reference values. MAE is in ppm and in
  units of the mean observational uncertainty
  $\bar\sigma_{\mathrm{obs}}=91.1$ ppm. Best value per column in bold.
  Both multi-fidelity rows were fitted with $n_L=10\,000$, verified against
  the run configuration and the worker-side LF sample check.}
  \label{tab:cv_metrics}
  \begin{tabular}{lcccccc}
    \toprule
      & \multicolumn{4}{c}{Mean per-wavelength NRMSE $(-)$}
      & \multicolumn{2}{c}{MAE} \\
    \cmidrule(lr){2-5}\cmidrule(lr){6-7}
    Architecture
      & all bins & F322W2 & F444W & MIRI
      & (ppm) & $/\bar\sigma_{\mathrm{obs}}$ \\
    \midrule
    \multicolumn{7}{@{}l}{\emph{Single fidelity (HF only)}}\\
    \quad per-$\lambda$ GP
      & 0.4953 & 0.5349 & 0.4758 & \textbf{0.3914} & 99.8 & 1.095 \\
    \quad per-$\lambda$ GPBoost
      & 0.4864 & 0.5282 & 0.4481 & 0.4192 & 93.4 & 1.025 \\
    \quad PCA--PCE$^{\dagger}$
      & 0.4826 & 0.5051 & 0.4704 & 0.4264 & 93.9 & 1.031 \\
    \addlinespace[2pt]
    \multicolumn{7}{@{}l}{\emph{Multi-fidelity}}\\
    \quad Model 1, per-$\lambda$ MFGP
      & 0.4837 & 0.5495 & \textbf{0.3933} & 0.4506 & 100.5 & 1.104 \\
    \quad Model 2, per-$\lambda$ GPBoost
      & \textbf{0.4350} & \textbf{0.4603} & 0.3956 & 0.4339 & \textbf{86.9} & \textbf{0.954} \\
    \bottomrule
  \end{tabular}

  \vspace{4pt}
  \begin{minipage}{0.95\textwidth}
    \footnotesize
    Instrument columns aggregate the $103$ bins of NIRCam F322W2
    ($2.46$--$3.99\,\mu$m), the $65$ bins of NIRCam F444W
    ($4.00$--$4.96\,\mu$m), and the $27$ bins of MIRI LRS
    ($5.25$--$11.75\,\mu$m). The all-bin averages weight bins equally;
    they do not weight by wavelength interval. The MIRI bin-centre spacing
    is $16.7$ times the NIRCam spacing. The Model~2 advantage over
    per-$\lambda$ GPBoost in the first column is not significant at the
    design level ($p=0.42$); see Section~\ref{sec:cv_results}.
    $^{\dagger}$Reference surrogate of \citet{Nussbaumer2026}, evaluated
    from the supplied leave-one-out predictions. Its wavelength headers match
    the HF grid and are checked in code; its row order is assumed to match the
    high-fidelity design order because the reference file contains no sample
    identifiers.
  \end{minipage}
\end{table}

\begin{figure}[t]
    \centering
    \includegraphics[width=\linewidth]{figures/cv_performance.pdf}
    \caption{Leave-one-out performance across the $195$ wavelength bins for
    all five architectures, over $97$ held-out HF designs. Upper panel:
    per-wavelength NRMSE, $\mathrm{RMSE}_j$ divided by the population standard
    deviation $s_j$ of the reference values at that wavelength. Lower panel:
    per-wavelength MAE in ppm. Solid curves are multi-fidelity models and the
    PCA--PCE reference; dashed curves are the two single-fidelity baselines.
    The NIRCam bands ($<5\,\mu$m) are densely sampled at $0.015\,\mu$m while
    MIRI ($\ge5.25\,\mu$m) is sampled at $0.25\,\mu$m, so the right-hand part
    of each panel is sparser. The two panels rank the models differently, and
    MAE rises steeply with wavelength while NRMSE does not, which is the
    amplitude effect quantified in Section~\ref{sec:results}.}
    \label{fig:nrmse_lambda}
\end{figure}

\resolved{results graphic}{The figure was located at
\texttt{results/cv/loo\_lf10000/manuscript/cv\_performance.pdf} (and
\texttt{.png}); copy it to \texttt{figures/cv\_performance.pdf} for the build.
Its contents are specified exactly by \texttt{report\_cv\_performance.make\_figure}:
two vertically stacked, $x$-shared panels of per-wavelength NRMSE and MAE in
ppm, one line per predictor, dashed for the single-fidelity rows. The caption
above describes that figure. The dangling \texttt{fig:nrmse\_lambda} label of
the previous draft refers to this same graphic --- it is the only
wavelength-resolved error figure produced --- so the label is attached here and
the placeholder float removed.}

\section{Use in Bayesian Retrieval}
\label{sec:bayesian-inversion}

Bayesian inversion combines a forward model with observations and prior
information to infer atmospheric parameters
\citep{stuart2010inverse, damien2013bayesian, ghanem2017handbook,
chiappetta2026efficient}. The simulator data $\mathcal D$ are defined in
\eqref{eq:inversion-training-data}; the astronomical measurement
$\mathbf y_{\mathrm{obs}}\in\mathbb R^m$ supplies the retrieval target.

The observation model is
\begin{equation}
    \mathbf y_{\mathrm{obs}}
    =
    \mathbf f_H(\bm\theta)
    +\bm\varepsilon_{\mathrm{obs}},
    \qquad
    \bm\varepsilon_{\mathrm{obs}}
    \sim
    \mathcal N_{m}
    \bigl(\mathbf 0,\bm\Sigma_{\mathrm{obs}}\bigr),
    \label{eq:inversion-observation-model}
\end{equation}
where $\mathbf f_H$ is the HF forward map and
$\bm\Sigma_{\mathrm{obs}}=\operatorname{diag}(\sigma^2_{\mathrm{obs},1},
\ldots,\sigma^2_{\mathrm{obs},m})$ is the observational error covariance.
This model assumes zero-mean, independent Gaussian measurement errors and
includes no additional discrepancy between the HF simulator and the physical
system. Bayes' theorem gives
\begin{equation}
    p(\bm\theta\mid\mathbf y_{\mathrm{obs}})
    \propto
    p(\mathbf y_{\mathrm{obs}}\mid\bm\theta)
    p(\bm\theta),
    \label{eq:inversion-exact-posterior}
\end{equation}
with $p(\bm\theta)$ the prior of Section~\ref{subsub:prior}; the evidence is
constant in $\bm\theta$ and is not evaluated when sampling by MCMC.

To reduce the cost of repeated forward evaluations, the HF simulator is
replaced by its surrogate predictive distribution conditional on $\mathcal D$,
\begin{equation}
    \mathbf f_H(\bm\theta)\mid\mathcal D
    \sim
    \mathcal N_{m}\!\left(
        \widehat{\mathbf f}_H(\bm\theta),
        \bm\Sigma_{\mathrm{em}}
    \right),
    \label{eq:inversion-emulator-distribution}
\end{equation}
where $\widehat{\mathbf f}_H(\bm\theta)$ is the surrogate predictive mean
spectrum and $\bm\Sigma_{\mathrm{em}}$ its error covariance. Assuming
observational and emulator errors independent, integration over the unknown HF
spectrum $\mathbf w$ yields
\begin{align}
    \widehat p(\mathbf y_{\mathrm{obs}}
        \mid\bm\theta,\mathcal D)
    &=
    \int_{\mathbb R^{m}}
    \mathcal N_{m}\!\left(
        \mathbf y_{\mathrm{obs}};
        \mathbf w,\bm\Sigma_{\mathrm{obs}}
    \right)
    \mathcal N_{m}\!\left(
        \mathbf w;
        \widehat{\mathbf f}_H(\bm\theta),
        \bm\Sigma_{\mathrm{em}}
    \right)d\mathbf w
    \nonumber\\
    &=
    \mathcal N_{m}\!\left(
        \mathbf y_{\mathrm{obs}};
        \widehat{\mathbf f}_H(\bm\theta),
        \bm\Sigma_{\mathrm{obs}}
        +\bm\Sigma_{\mathrm{em}}
    \right).
    \label{eq:inversion-surrogate-likelihood}
\end{align}

\subsection{The implemented emulator error covariance}
\label{subsec:implemented-likelihood}

Section~\ref{sec:calibration} showed that the GP predictive variances are not
calibrated, so $\bm\Sigma_{\mathrm{em}}$ is not taken from
\eqref{eq:model1-var} or \eqref{eq:model2-var}. It is instead estimated from
the $97\times195$ matrix of leave-one-out residuals $\{e_{ij}\}$ of the same
model family that supplies the forward mean. These are genuine out-of-sample
errors, so the error budget is measured rather than assumed. Because $97$
samples cannot support a $195\times195$ sample covariance, the estimate uses
Ledoit--Wolf shrinkage, about the residual mean rather than about zero, since
the emulator is biased as well as noisy:
\begin{equation}
    \widehat{\bm\Sigma}_{\mathrm{em}}
    =\mathrm{LedoitWolf}\bigl(\{ \mathbf e_i \}_{i=1}^{97}\bigr),
    \qquad
    \bm\Sigma=\operatorname{diag}(\sigma^2_{\mathrm{obs},j})
    +\widehat{\bm\Sigma}_{\mathrm{em}}.
    \label{eq:sigma-em}
\end{equation}
Three properties distinguish this from the construction proposed in the
previous draft. It is a \emph{full} matrix, not diagonal: emulator residuals
are coherent across wavelength, and treating $195$ correlated errors as
independent would shrink the posterior spuriously. It is \emph{independent of}
$\bm\theta$, so the likelihood does not factorise over wavelengths, but its
Cholesky factor and log-determinant are computed once and reused, making each
likelihood evaluation one triangular solve. And it is \emph{empirical},
sidestepping the calibration failure rather than inheriting it. A diagonal
variant is retained in the code as a sensitivity check.

Because $\bm\Sigma$ does not depend on $\bm\theta$, its log-determinant is an
additive constant and does not affect the posterior shape --- a simplification
the $\bm\theta$-dependent construction would not have permitted. The
implementation nonetheless includes the full Gaussian normalisation.

\resolved{validation and final fitting; retrieval status}{Three corrections.
First, the previous draft presented the retrieval as proposed and unevaluated;
retrievals have in fact been run, and are reported in
Section~\ref{sec:retrieval-results} ---
\texttt{results/bayesian\_inversion/} holds completed chains for a synthetic
target and the observed spectrum, for both model families. Second, the
implemented likelihood is \eqref{eq:sigma-em}, not the diagonal
$\bm\theta$-dependent $\bm\Sigma_{\mathrm{sur}}(\bm\theta)$ of the draft; the
per-wavelength factorisation \eqref{eq:model1-wavelength-likelihood} therefore
does \emph{not} apply to the runs reported here, and is retained below only as
the special case it is. Third, the draft's instruction to validate the
predictive variance before treating it as calibrated has been carried out in
Section~\ref{sec:calibration}, with a negative result, which is precisely why
the empirical covariance is used. The draft's caution about retaining the
log-determinant is correct in general but moot for a $\bm\theta$-independent
covariance.}

\subsection{Prior distributions}
\label{subsub:prior}

The independent priors are specified in the coordinates of
Table~\ref{tab:inputs}, which are the coordinates of the training design, so
no conversion is applied:
\[
    \bm\theta =
    \bigl(
        \log_{10}K_{zz}, R_P, T_{\mathrm{int}}, \mathrm{C/O},
        [\mathrm{N/H}], [\mathrm{O/H}], [\mathrm{S/H}], \log_{10}g, \fhr
    \bigr)^\top,
    \qquad
    p(\bm\theta)=\prod_{q=1}^{d}p_q(\theta_q).
\]
The marginal distributions are
\begin{align}
    \log_{10}K_{zz} &\sim \mathcal U(8,11.5), &
    T_{\mathrm{int}} &\sim \mathcal U(150,450), \\
    \mathrm{C/O} &\sim \mathcal U(0.1,1.5), &
    \fhr &\sim \mathcal U\!\left(\tfrac14,\tfrac23\right), \\
    [\mathrm{N/H}],\;[\mathrm{O/H}],\;[\mathrm{S/H}]
    &\overset{\mathrm{ind}}{\sim} \mathcal U(-1,2), \\
    R_P &\sim \mathcal{TN}_{[0.882,\,1.120]}(0.999,0.031^2), &
    \log_{10}g &\sim \mathcal{TN}_{[3.085,\,3.208]}(3.145,0.016^2).
\end{align}
Here $\mathcal{TN}_{[a,b]}(\mu,\sigma^2)$ denotes a
$\mathcal N(\mu,\sigma^2)$ distribution truncated and renormalised to $[a,b]$,
and $0.031$ and $0.016$ are the standard deviations before truncation.
$R_P$ is in Jupiter radii, $T_{\mathrm{int}}$ in kelvin, and the logarithms of
$K_{zz}$ and $g$ refer to values in $\mathrm{cm^2\,s^{-1}}$ and
$\mathrm{cm\,s^{-2}}$. Elemental abundances use the solar-relative logarithmic
convention
\[
    [\mathrm{X/H}]
    =
    \log_{10}\!\left(
        \frac{n_{\mathrm X}/n_{\mathrm H}}
             {(n_{\mathrm X}/n_{\mathrm H})_\odot}
    \right),
\]
so the carbon abundance is determined jointly by $\mathrm{C/O}$ and
$[\mathrm{O/H}]$. Prior independence is imposed in these coordinates.

The truncation bounds on $R_P$ and $\log_{10}g$ are the \emph{emulator
support}, not the nominal astrophysical bounds. The surrogate is trained only
inside the LF design box, and extrapolates meaninglessly outside it, so every
proposal is restricted to that box. For the seven uniform parameters the box
and the prior bounds coincide exactly. For the two Gaussians the box is the
LF design range, $\pm3.85$ standard deviations on both, so truncating there
discards about $10^{-4}$ of the prior mass.

\resolved{coordinates, prior provenance and support}{Read from
\texttt{bayesian\_inversion/priors.py} (\texttt{PRIOR\_SPEC} and
\texttt{EMULATOR\_BOX}). The previous draft stated truncation intervals
$R_P\in[0.1,2]$ and $\log_{10}g\in[2,6]$; those are the nominal spec, whereas
the implemented sampler truncates to $[0.881978, 1.120004]$ and
$[3.084744, 3.207999]$ respectively, and rejects any draw outside the box. The
draft's concern about a mismatch between retrieval and training coordinates
does not arise: they are identical (Table~\ref{tab:inputs}), and
\texttt{priors.check\_box\_against\_design} tests the prior against the design
file directly. These are the priors actually used in the runs reported below,
not proposals.}

\subsection{Assembling the predictive mean}
\label{subsubsec:model1-posterior}

At each candidate $\bm\theta$, all $m=195$ wavelength emulators are evaluated
at the same atmospheric input and their predictive means assembled into
\begin{equation}
    \widehat{\mathbf f}_H(\bm\theta)
    =
    \bigl(
        \widehat f_{H,1}(\bm\theta),\ldots,
        \widehat f_{H,m}(\bm\theta)
    \bigr)^\top,
\end{equation}
which enters \eqref{eq:inversion-surrogate-likelihood} with the covariance
\eqref{eq:sigma-em}. Only the means are taken from the emulators; their
predictive variances are not used, for the reason given in
Section~\ref{sec:calibration}.

If one were instead to use the models' own predictive variances and treat the
bins as independent --- the construction of the previous draft --- the
likelihood would factorise as
\begin{equation}
    \widehat p(\mathbf y_{\mathrm{obs}}
        \mid\bm\theta,\mathcal D)
    =
    \prod_{j=1}^{m}
    \mathcal N\!\left(
        y_{\mathrm{obs},j};
        \widehat f_{H,j}(\bm\theta),
        \sigma_{\mathrm{obs},j}^{2}
        +\widehat\sigma^{2}_{H,j}(\bm\theta)
    \right).
    \label{eq:model1-wavelength-likelihood}
\end{equation}
This is not used here. It would discard the cross-wavelength correlation of
the emulator residuals, and it would inherit the miscalibration measured in
Section~\ref{sec:calibration}, in the anticonservative direction.

\subsection{Retrieval results}
\label{sec:retrieval-results}

The forward emulator for retrieval is the Model~2 multi-fidelity fit trained
on $10\,000$ LF rows and $96$ HF cases, holding out HF design point $81$ so
that it can serve as a closure test. Sampling uses \texttt{emcee} with $48$
walkers.

At the synthetic target, emulator error dominates the budget: the emulator's
out-of-sample RMS error at that design is $1.06\times10^{-4}$ against a median
observational error of $8.2\times10^{-5}$, so omitting
$\widehat{\bm\Sigma}_{\mathrm{em}}$ would tell the sampler the emulator is
exact and let it chase emulator artefacts. The closure posterior covers the
true parameters in $8$ of $9$ marginals at $95\%$ and $7$ of $9$ at $68\%$.
The exception is $\log_{10}K_{zz}$, whose true value $11.34$ lies just above
the $95\%$ upper limit $11.23$, near the edge of the prior box. Posterior
widths relative to the prior range vary widely: $R_P$ and $\log_{10}g$ are
prior-dominated by construction, $\mathrm{C/O}$ and $[\mathrm{O/H}]$ narrow to
roughly $0.38$--$0.39$ of the prior range, while $T_{\mathrm{int}}$,
$[\mathrm{N/H}]$ and $\log_{10}K_{zz}$ remain close to their priors, at
$0.88$--$0.93$.

At the observed spectrum, the same emulator yields a posterior median
$\mathrm{C/O}=0.54$ with $95\%$ interval $[0.25,0.85]$, and
$[\mathrm{O/H}]=0.80$ with $[0.39,1.59]$; $T_{\mathrm{int}}$,
$[\mathrm{N/H}]$, $[\mathrm{S/H}]$ and $\log_{10}K_{zz}$ remain
prior-dominated, retaining $0.69$--$0.92$ of the prior width.

\paragraph{Convergence.} These posteriors are \emph{not} converged and must be
read as provisional. The longest integrated autocorrelation time is $949$
steps against a post-burn-in chain of $11\,963$ steps, i.e. $13\tau$, where
\texttt{emcee}'s rule of thumb asks for $50\tau$; the run configuration records
\texttt{converged: false}. An autocorrelation time estimated from a chain this
short is itself biased low, so $13\tau$ is an optimistic figure. The intervals
above should not be quoted as final credible intervals, and no scientific
conclusion about WASP-80b's atmosphere is drawn from them here. Longer chains
are the obvious next step.

\section{Discussion}
\label{sec:disc}

The comparison shows that a boosted mean improves spectral prediction, and
that the evidence for fidelity coupling specifically is weaker than the
headline table suggests. Model~2 has the lowest all-bin mean NRMSE and the
lowest MAE, and beats Model~1 decisively on paired design scores. But it does
not separate from the HF-only GPBoost baseline at the design level
($p=0.42$), and Model~1 MF does not separate from its own single-fidelity
counterpart ($p=0.12$). Read together with the aggregate ranking, the most
defensible statement is that the mean model matters and the fidelity coupling,
at this LF budget and with this kernel, does not yet demonstrably help. Since
the LF budget is already at its maximum of $10\,000$ rows, the remaining
headroom is structural rather than a matter of more LF data.

The aggregate ranking also depends on the error summary. Equal-bin averaging
gives most weight to NIRCam, whereas errors in ppm are strongly influenced by
longer-wavelength amplitudes; MAE and NRMSE order the architectures
differently, and the band-wise minima fall to three different models. The
paired design-level comparison does not establish a clear advantage over
PCA--PCE. These qualifications limit the interpretation of the lowest all-bin
average to the particular data, metrics and protocol used here.

The uncertainty results are a clearer negative. A mean absolute error of
$86.9$ ppm is substantial relative to the mean measurement uncertainty of
$91.1$ ppm, so emulator error cannot be ignored in a retrieval; but the models'
own predictive variances are anticonservative in the tails, and the boosted
models occasionally report near-zero variance. Emulator error had to be
measured from held-out residuals instead. The likelihood used in the retrieval
carries cross-wavelength emulator correlation but omits, still, uncertainty in
the fitted mean and covariance parameters and any discrepancy between the HF
simulator and the physical atmosphere. It also treats emulator error as
homoscedastic in $\bm\theta$, which Section~\ref{sec:results} shows is
optimistic in the high-$[\mathrm{O/H}]$ region.

Two structural limitations follow from the modelling choices rather than the
data. The kernel is isotropic Matérn-$3/2$ with one range for all nine inputs,
so the surrogate cannot express differing sensitivity across parameters; an
ARD variant is available in the same software and is the obvious first
experiment. And both models fit $195$ independent scalar GPs, so nothing
enforces smoothness along wavelength --- which is visible in the residuals and
is what Appendix~\ref{subsec:model_augmented} addresses.

The interpretation of generalisation depends on data provenance. The LF
spectra come from profile emulators trained on HF cases, and those emulators
are not rebuilt within the folds, so the multi-fidelity rows measure
generalisation conditional on that pretrained resource. The evidence limiting
this concern is that no HF input appears in the LF design and that the
single-fidelity rows are unaffected. The PCA--PCE comparison additionally
depends on an assumed row alignment, since the reference file carries no
sample identifiers.

\section{Conclusion and Outlook}

This study compares per-wavelength autoregressive and boosted multi-fidelity
emulators for WASP-80b emission spectra. The boosted MF model attains a mean
per-wavelength NRMSE of $0.4350$ and an MAE of $86.9$ ppm, the best of the five
architectures on both, but its advantage over the HF-only boosted baseline is
not significant at the design level, and gains vary across spectral bands. The
improvement is attributable to the boosted mean rather than demonstrably to
fidelity coupling.

The predictive variances of all four fitted models are miscalibrated, so the
retrieval likelihood carries an emulator error covariance estimated from
leave-one-out residuals, retaining cross-wavelength correlation. Retrievals of
a synthetic and the observed spectrum have been run under this likelihood; the
synthetic closure test covers $8$ of $9$ true parameters at $95\%$, but the
chains are not converged and the intervals are provisional.

The immediate next steps are to run the chains to the $50\tau$ standard, to
confirm the upstream LF provenance so the generalisation claim can be stated
without qualification, and to test an ARD kernel and joint rather than staged
estimation of the boosted model. The wavelength-augmented construction remains
an extension requiring its own spectral holdout experiment.

\bibliography{bibliography.bib}
\bibliographystyle{abbrvnat}
\appendix

\section{Scalar autoregressive Gaussian process background}
\label{app:ar_background}
\label{subsec:ar_gp}

This appendix records the scalar autoregressive construction underlying the
per-wavelength models. The formulation follows \citet{KennedyOHagan2000};
recursive co-kriging is developed by \citet{LeGratietGarnier2014} under design
and model assumptions that, as noted in Section~\ref{sec:data}, do not hold
here. Let $\bm\theta\in\mathcal X\subset\mathbb R^d$ be the atmospheric input
vector and consider one scalar response at each fidelity. The autoregressive
model expresses the HF process as a scaled LF process plus an independent
discrepancy,
\begin{equation}
    f_H(\bm\theta) = \rho\, f_L(\bm\theta) + \delta(\bm\theta),
    \label{eq:ar1}
\end{equation}
where $\rho\in\mathbb{R}$ is a scalar coupling coefficient and
$\delta(\bm\theta)$ absorbs the systematic bias of the LF simulator. The
recursive form of the construction rests on the Markov assumption of
\citet{LeGratietGarnier2014},
\begin{equation}
    \operatorname{Cov}\!\big(f_H(\bm\theta),\, f_L(\bm\theta')
    \,\big|\, f_L(\bm\theta)\big) = 0,
    \qquad \forall\,\bm\theta'\neq\bm\theta,
    \label{eq:markov}
\end{equation}
i.e.\ once $f_L(\bm\theta)$ is known, LF evaluations elsewhere carry no further
information about $f_H(\bm\theta)$. The models in this paper are fitted as the
joint two-fidelity GP below, which does not require \eqref{eq:markov} or nested
designs. The two components are assigned independent Gaussian process priors,
\begin{align}
    f_L    &\sim \mathcal{GP}\!\big(\mu_L,\,
              k_L(\bm\theta,\bm\theta';\,\bm\psi_L)\big),
              \label{eq:prior_L}\\
    \delta &\sim \mathcal{GP}\!\big(\mu_\delta,\,
              k_\delta(\bm\theta,\bm\theta';\,\bm\psi_\delta)\big),
              \qquad \delta \perp f_L.
              \label{eq:prior_delta}
\end{align}
A general stationary covariance may be written
\begin{equation}
    k(\bm\theta,\bm\theta')
    = \sigma^{2}\,\varphi\!\left(
      \left[\textstyle\sum_{q=1}^{d}
      (\theta_q-\theta_q')^{2}/\ell_q^{2}\right]^{1/2}\right),
    \label{eq:general-kernel}
\end{equation}
with $\varphi$ the correlation function, $\sigma^{2}$ the signal variance and
$\{\ell_q\}$ per-dimension length-scales. Two specialisations matter here.
Taking $\varphi(r)=\exp(-r^2/2)$ with distinct $\ell_q$ gives the
squared-exponential ARD kernel; taking
$\varphi(r)=(1+\sqrt3 r)\exp(-\sqrt3 r)$ with a common
$\ell_q\equiv\ell$ gives the isotropic Matérn-$3/2$ kernel
\eqref{eq:matern32}, which is the one used in every reported run. Substituting
\eqref{eq:prior_L}--\eqref{eq:prior_delta} into \eqref{eq:ar1}, the joint
process $(f_L,f_H)^{\top}$ is Gaussian,
\begin{equation}
    \begin{bmatrix} f_L(\bm\theta)\\[2pt] f_H(\bm\theta)\end{bmatrix}
    \sim \mathcal{GP}\!\left(
    \begin{bmatrix} \mu_L \\[2pt] \rho\,\mu_L + \mu_\delta \end{bmatrix},\;
    \begin{bmatrix}
        k_L(\bm\theta,\bm\theta') & \rho\,k_L(\bm\theta,\bm\theta') \\[4pt]
        \rho\,k_L(\bm\theta,\bm\theta') &
        \rho^{2}\,k_L(\bm\theta,\bm\theta') + k_\delta(\bm\theta,\bm\theta')
    \end{bmatrix}
    \right),
    \label{eq:joint}
\end{equation}
so that the cross-covariance between fidelities is $\rho\,k_L$, the HF marginal
variance is $\rho^{2}k_L+k_\delta$, and $\rho$ specifies the multiplicative
coupling. The amount of information transferred also depends on the covariance
scales and data.

\paragraph{Fitted couplings.} The estimated $\rho_j$ differ markedly between
the two models. For Model~1 MF the median across the $97\times195$ fits is
$1.135$ (interquartile range $0.854$--$1.851$, full range $0.428$--$2.774$),
never negative; for Model~2 MF the median is $0.815$ (interquartile range
$0.590$--$0.987$, range $-0.349$--$1.634$), with $3.1\%$ of fits negative.
Model~1 must use $\rho$ to reconcile the LF and HF amplitude difference
directly, which pushes it above one; in Model~2 the shared boosted mean
already absorbs that offset, leaving $\rho$ to describe residual coupling and
concentrating it below one. This is the mechanism by which the mean model and
the fidelity coupling interact.

\section{Wavelength-augmented extension}
\label{subsec:model_augmented}

This appendix presents a formulation that is \emph{not} among the evaluated
models: it has no reported results, and no gap-filling experiment was run.

Both main-text models treat wavelength bins independently. An alternative
is to represent the spectrum as a function $s_t(\bm\theta;\lambda)$ on
$\Lambda=[\lambda_{\min},\lambda_{\max}]$, here approximately
$[2.46,11.75]\,\mu\mathrm{m}$, and model wavelength dependence explicitly.
This introduces a smoothness assumption along wavelength through the
kernel. Define the augmented input
\begin{equation}
    \mathbf{z} = (\bm\theta,\lambda) \in \mathcal{Z}
    = \mathcal{X}\times\Lambda \subset \mathbb{R}^{d+1},
    \label{eq:model2_augmented_input}
\end{equation}
and view each fidelity as a single \emph{scalar}-valued map on
$\mathcal{Z}$, $f_t(\mathbf{z}) = s_t(\bm\theta;\lambda)$, $t\in\{L,H\}$. The
autoregressive construction of Appendix~\ref{subsec:ar_gp} is then applied
once, on the augmented domain:
\begin{equation}
    f_{H}(\bm\theta,\lambda)
    = \rho\, f_{L}(\bm\theta,\lambda) + \delta(\bm\theta,\lambda),
    \qquad
    f_{L} \sim \mathcal{GP}(0,k_{L}),
    \quad
    \delta \sim \mathcal{GP}(0,k_{\delta}),
    \quad
    \delta \perp f_{L},
    \label{eq:model2_ar}
\end{equation}
with the covariance \eqref{eq:general-kernel} now taken over $d+1$ dimensions,
so that each kernel acquires one additional length-scale along the spectral
axis. The coupling $\rho$ is a single scalar for the entire spectrum; any
wavelength dependence of the LF$\to$HF transfer that a scalar $\rho$ cannot
express is absorbed by $\delta(\bm\theta,\lambda)$, which is free to vary along
$\lambda$.

The parameter accounting depends on which kernel is used, and the previous
draft's figures assumed the ARD form. Under an ARD kernel with two
fidelity-specific noise variances, the augmented model carries
$2\times(d+2)+1=2\times11+1$ covariance parameters plus two noise variances,
i.e.\ $25$, against Model~1's bank of $m(2d+5)=195\times23=4485$. Under the
kernel actually used in this paper --- isotropic Matérn-$3/2$ with one shared
noise variance --- the corresponding counts are $6$ for the augmented model
against $195\times6=1170$ for the bank, using \eqref{eq:implemented-params}.
Either way the reduction is by a factor of roughly $180$--$195$, and the
augmented model induces non-zero covariance across bins,
\begin{equation}
    \operatorname{Cov}\!\bigl(f_H(\bm\theta,\lambda_j),\,
    f_H(\bm\theta',\lambda_{j'})\bigr)
    = \rho^{2}k_{L}\bigl((\bm\theta,\lambda_j),(\bm\theta',\lambda_{j'})\bigr)
    + k_{\delta}\bigl((\bm\theta,\lambda_j),(\bm\theta',\lambda_{j'})\bigr),
    \label{eq:model2_crosscov}
\end{equation}
so the model can share information across wavelengths. The parameter count
refers to covariance, coupling and noise parameters, not to the number of
observations or the computational cost: using all $10\,000$ LF and $97$ HF
spectra would yield $n=(10\,000+97)\times195=1\,968\,915$ scalar observations.

\openitem{bin representation}{Whether the targets are point evaluations at bin
centres or averages over the instrument bandpass is not recorded in the data
files, and it matters here specifically: interpolating along $\lambda$ through
a kernel presumes a pointwise response, and bin widths differ by a factor of
$16.7$ between NIRCam and MIRI. This must be settled before the augmented model
is fitted, though not before the main-text models, which never interpolate
along wavelength.}

Both fidelities are represented on the complete Cartesian product of
their atmospheric design with the full wavelength grid
$\Lambda_m=\{\lambda_1,\ldots,\lambda_m\}$,
\begin{align}
\label{eq:model2_designs_full}
    \mathbf{Z}_{L} = \Theta_{L}\times\Lambda_m, \qquad
    \mathbf{Z}_{H} = \Theta_{H}\times\Lambda_m,
\end{align}
so that every sample--wavelength pair contributes one augmented row, with
$|\mathbf{Z}_L| = n_L\,m$ and $|\mathbf{Z}_H| = n_H\,m$. A binary fidelity
indicator is appended as an additional column; it selects covariance blocks in
\eqref{eq:model2_blockcov} and is not a continuous kernel coordinate. The
$n = (n_L + n_H)\,m$ targets are stacked LF above HF into a single vector
$\mathbf{y}\in\mathbb{R}^{\,n}$. Unlike the main-text models, which leave the
response untransformed, this construction would need one global centring and
scaling before fitting, because the response amplitude varies by a factor of
$22$ along the wavelength axis that is now a kernel coordinate; predictions
would be mapped back afterwards.

Under \eqref{eq:model2_ar} plus independent Gaussian observation noise with
fidelity-specific variances, the stacked vector is Gaussian,
$\mathbf{y}\sim\mathcal{N}(\mathbf 0,\bm\Sigma(\rho,\bm\eta))$, with the same
two-fidelity block structure as \eqref{eq:model1_blockcov}, now evaluated on
the augmented designs:
\begin{equation}
    \bm\Sigma =
    \begin{bmatrix}
        \mathbf{K}_{L}^{LL} + \sigma^{2}_{\varepsilon,L}\mathbf{I}
        & \rho\,\mathbf{K}_{L}^{LH} \\[4pt]
        \rho\,\mathbf{K}_{L}^{HL}
        & \rho^{2}\,\mathbf{K}_{L}^{HH}
          + \mathbf{K}_{\delta}^{HH}
          + \sigma^{2}_{\varepsilon,H}\mathbf{I}
    \end{bmatrix},
    \label{eq:model2_blockcov}
\end{equation}
where $[\mathbf{K}_{L}^{AB}]_{ii'}=k_L(\mathbf z_i^A,\mathbf z_{i'}^B)$. The
kernels, $\rho$ and the noise variances would be estimated jointly by
gradient-based optimisation of the single log marginal likelihood
\begin{equation}
    \log L_{\mathrm{aug}}(\rho,\bm\eta)
    = -\tfrac{1}{2}\,\mathbf{y}^{\top}\bm\Sigma^{-1}\mathbf{y}
      -\tfrac{1}{2}\log\det\bm\Sigma
      -\tfrac{n}{2}\log 2\pi .
    \label{eq:model2_lml}
\end{equation}
Writing $b,b'\in\{0,1\}$ for the fidelity indicators of two rows, the joint
kernel's LF component is $\rho^{\,b+b'}k_{L}(\mathbf{z},\mathbf{z}')$. If
optimisation restricts $\rho>0$ and uses $\log\rho$, then
\begin{equation}
    \frac{\partial k_{\mathrm{MF}}\bigl((\mathbf{z},b),(\mathbf{z}',b')\bigr)}{\partial \log\rho}
    = (b+b')\,\rho^{\,b+b'}\,k_{L}(\mathbf{z},\mathbf{z}'),
    \label{eq:model2_rho_grad}
\end{equation}
i.e.\ the LL block contributes no $\rho$-gradient, the cross blocks scale
linearly, and the HH block quadratically, matching
\eqref{eq:model2_blockcov}. Note that the GPBoost implementation used for the
main-text models leaves $\rho$ unrestricted in sign, and $3.1\%$ of Model~2's
fitted couplings are negative, so a $\log\rho$ parameterisation would not be
appropriate without a change of model.

For prediction, a new atmospheric input $\bm\theta^{*}$ is expanded to the full
wavelength grid, $\{(\bm\theta^{*},\lambda_{j})\}_{j=1}^{m}$, and the HF
posterior follows from Gaussian conditioning of $f_H$ on $\mathbf{y}$ under
\eqref{eq:model2_blockcov}.

\openitem{implementation and evaluation if pursued}{At $n\approx2\times10^{6}$
rows this model is not fittable exactly; a Vecchia approximation of the kind
already used here, or a training subset, would be required, and the actual
scaling rule must be documented rather than assumed. Evaluation must hold out
wavelength blocks or a specified missing-bin pattern to demonstrate spectral
interpolation, and must additionally hold out atmospheric inputs if both forms
of generalisation are claimed. Neither a fitting procedure nor successful gap
filling is established by this formulation alone.}

\end{document}
